%!TEX program = lualatex
%!TEX encoding = UTF-8 Unicode  
%!TEX root = chapter13.tex  
\documentclass[main.tex]{subfiles}
\begin{document} 
%----------------------------------------------------------------------------------------
%	CHAPTER 13
%----------------------------------------------------------------------------------------
\cleardoublepage
\loadgeometry{marginlayout}  
\pagestyle{scrheadings} % Print headers again  
\KOMAoptions{
headwidth=textwithmarginpar,
headsepline=true,
headsepline= 0.5pt,
footwidth=textwithmarginpar,
}   
\onehalfspacing 
\setcounter{chapter}{12}



\chapter{Vecteurs (2) Opérations et colinéarité}
 
\section{Produit d'un vecteur par un réel}

\begin{definition}[interprétation géométrique] $\vv{u}\neq \vv{0}$. 
Le produit d'un réel $k$ par un vecteur $\vv{u}$ est un vecteur noté $k\vv{u}$ : 
\begin{description}
\item[P1]  $0\vv{u}=\vv{0}$.
\item[P2]  Si $k \neq  0$,  \og $k\vv{u}$ \fg{} désigne le vecteur :
\begin{enumerate}[label=\textbullet, leftmargin=*]
\item ayant même direction que $\vv{u}$ 
\item $\norm{k\vv{u}}=\abs{k}\times \norm{\vv{u}}$
\item Si $k>0$ alors $k\vv{u}$ et $\vv{u}$ ont \textbf{même sens}.\\
Si $k<0$ alors $k\vv{u}$ et $\vv{u}$ sont de \textbf{sens contraires}
\end{enumerate}
\end{description} 
\end{definition}


\vfill

\begin{definition}[vecteurs opposés]  $\vv{v}=-\vv{u}$ signifie que les vecteurs $\vv{u}$ et $\vv{v}$  ont  même direction ,  même norme mais sont de sens contraires. 
En particulier : 	$-\vv{AB}=\vv{BA}$
\end{definition} 
\begin{example} $I$ est le milieu du segment $[AB]$. En utilisant uniquement des vecteurs d'extrémités  $A$, $B$ ou $I$ :
\begin{enumerate}[label=\alph*), leftmargin=*]
\item Énumérer 2 égalités de vecteurs.
\vfill
\item Énumérer 3 paires de vecteurs opposés $\vv{v}=-\vv{u}$
\vfill
\item Énumérer 3 paires de vecteurs qui peuvent s'écrire $\vv{v}=2\vv{u}$
\vfill
\item Énumérer 3 paires de vecteurs qui peuvent s'écrire $\vv{v}=\frac{1}{2}\vv{u}$
\vfill
\end{enumerate}
\end{example}
\vfill 
\begin{tBox} Soit un repère $(O;I,J)$, et le vecteur $\vv{u}\begin{pmatrix}
x\\y
\end{pmatrix}$.

\noindent	Le vecteur $k\vv{u}$ a pour coordonnées $k\vv{u}\begin{pmatrix}
kx\\ky
\end{pmatrix}$. 
%\noindent On pose $\vec{\imath}=\vv{OI}\begin{pmatrix}
%	1\\0
%\end{pmatrix}$ et $\vec{\jmath}=\vv{OJ}\begin{pmatrix}
%0\\1
%\end{pmatrix}$. 
%
%\noindent Le couple de vecteurs $(\vec{\imath}, \vec{\jmath}=\vv{OJ})$ est une base. Le vecteur $\vv{u}\begin{pmatrix}
%	x\\y
%\end{pmatrix}$ s'écrit comme la combinaison linéaire  $\vv{u}=x\vec{\imath}+y\vec{\jmath}$.
\end{tBox}   
%----------------------------------------------------------------------------------------
%	 
%----------------------------------------------------------------------------------------

\newpage 
\loadgeometry{widelayout}  
\pagestyle{scrheadings} % Print headers again  
\KOMAoptions{
headwidth=textwithmarginpar,
headsepline=true,
headsepline= 0.5pt,
footwidth=textwithmarginpar,
} 
\onehalfspacing


\subsection*{Exercices : produit d'un vecteur par un réel}
\begin{example}\hfill \\
Les vecteurs ci-dessous ont tous la même direction.\\
\begin{tikzpicture}[scale=1]
\tkzInit[xmax=6,ymax=1.5]%\tkzGrid[sub]
\tkzSetUpPoint[size=4, fill=black]
\tkzfctset{tan style/.style={->,>=latex, color=red, line width=1.2pt}}  

\def \xa{3}; 
\def \ya{0}; 

\def \xb{5};
\def \yb{0};

\tkzDefPoint(1,0){A} \tkzDefShiftPoint[A](\xa,\ya){B}
\tkzDefPoint(0,-1){C} \tkzDefShiftPoint[C](\xb,\yb){D}

\tkzDrawSegments[->, >=latex, line width=1.2pt](A,B C,D)


\tkzLabelPoints[below, font=\scriptsize](C,D) 
\tkzLabelPoints[above, font=\scriptsize](A,B) 

\foreach \i in {1,2}{%
\tkzMarkSegment[pos=\i/3,mark=|,  size=2pt](A,B)
}

\foreach \i in {1,...,4}{%
\tkzMarkSegment[pos=\i/5,mark=|,  size=2pt](C,D)
} 
\end{tikzpicture} 
\hfill 
\begin{tikzpicture}[scale=1]
\tkzInit[xmax=6,ymax=1.5]%\tkzGrid[sub]
\tkzSetUpPoint[size=4, fill=black]
\tkzfctset{tan style/.style={->,>=latex, color=red, line width=1.2pt}}  

\def \xa{3.6}; 
\def \ya{0}; 

\def \xb{2.4};
\def \yb{0};

\tkzDefPoint(0,0){A} \tkzDefShiftPoint[A](\xa,\ya){B}
\tkzDefPoint(3,-1){C} \tkzDefShiftPoint[C](-\xb,\yb){D}
\tkzDrawSegments[->, >=latex, line width=1.2pt](A,B C,D)

\tkzLabelPoints[below, font=\scriptsize](C,D) 
\tkzLabelPoints[above, font=\scriptsize](A,B) 

\foreach \i in {1,2}{%
\tkzMarkSegment[pos=\i/3,mark=|,  size=2pt](A,B)
}

\foreach \i in {1}{%
\tkzMarkSegment[pos=\i/2,mark=|,  size=2pt](C,D)
} 
\end{tikzpicture} 
\hfill 
\begin{tikzpicture}[scale=1]
\tkzInit[xmax=6,ymax=1.5]%\tkzGrid[sub]
\tkzSetUpPoint[size=4, fill=black]
\tkzfctset{tan style/.style={->,>=latex, color=red, line width=1.2pt}}  

\def \xa{3}; 
\def \ya{0}; 

\def \xb{3};
\def \yb{0};

\tkzDefPoint(0,0){A} \tkzDefShiftPoint[A](\xa,\ya){B}
\tkzDefPoint(2,-1){C} \tkzDefShiftPoint[C](-\xb,\yb){D}
\tkzDrawSegments[->, >=latex, line width=1.2pt](A,B C,D)
\tkzDrawSegments[dashed, line width=0.5pt](A,D C,B)
\tkzLabelPoints[below, font=\scriptsize](C,D) 
\tkzLabelPoints[above, font=\scriptsize](A,B) 
\end{tikzpicture}  
\begin{minipage}{0.3\columnwidth} 
\begin{align*}
\frac{CD}{AB}&=\frac{5}{3} \\
CD & = \tfrac{5}{3}AB \\
\vv{CD}& =\tfrac{5}{3}\vv{AB}   
\end{align*}
\end{minipage}
\end{example}

\begin{exercise}\hfill 
\begin{enumerate}[label=\alph*),leftmargin=*]
\item Sur la figure ci-dessous : 
\begin{minipage}{0.8\columnwidth}%
\centering
\begin{tikzpicture} 
\tkzInit[xmin=-3.5, xmax=6.5]
\tkzDrawX[label={}, line width=1.2pt] %noticks, 
\tkzSetUpPoint[size=7, fill=black]
\tkzfctset{tan style/.style={->,>=latex, color=red, line width=1.2pt}} 
\foreach \x/\xtext in {-3/A, -1/B, 6/C} {
\draw (\x,0.1cm) -- (\x,-0.1cm) node[above] {$\xtext\strut$};
} 
\end{tikzpicture}
\end{minipage} 
%	
\begin{align*}
{AC}&= \ldots \ldots  {BC}
& {BC}&=\ldots \ldots  {AC}
&  {BC}&=\ldots \ldots  {BA} \\
\vv{AC}&= \ldots \ldots \vv{BC}
&\vv{BC}&=\ldots \ldots \vv{AC}
& \vv{BC}&=\ldots \ldots \vv{BA} \\
{CA}&=\ldots \ldots {BA}
& {AB}&=\ldots \ldots {BC} 
& {BA}&=\ldots \ldots {BC} \\
\vv{CA}&=\ldots \ldots \vv{BA}
& \vv{AB}&=\ldots \ldots \vv{BC} 
& \vv{BA}&=\ldots \ldots \vv{BC}  
\end{align*}  
\item À partir de la figure ci-dessous, compléter les égalités vectorielles suivantes:
\\
\begin{minipage}{1\columnwidth}%
\centering
\begin{tikzpicture} 
\tkzInit[xmin=-3.5, xmax=6.5]
\tkzDrawX[label={}, line width=1.2pt] %noticks, 
\foreach \x/\xtext in {-2/A, 1/B, 6/C} {
\draw[thick] (\x,0.1cm) -- (\x,-0.1cm) node[above] {$\xtext\strut$};
}
\end{tikzpicture}
\end{minipage} 
\begin{align*}
\vv{AC}& = \ldots \ldots \vv{BC}
&\vv{BC}& = \ldots \ldots \vv{BA}
& \vv{CA}& = \ldots \ldots \vv{AB}\\
% 		\\
\vv{BC}& = \ldots \ldots \vv{AC}
& \vv{BA}& = \ldots \ldots \vv{BC}
& \vv{BA}& = \ldots \ldots \vv{AB}
\end{align*} 
\item Le point $B$ est le milieu de $[AC]$. Compléter les égalités vectorielles suivantes :
\begin{align*}
\vv{BA}& = \ldots \ldots \vv{AC}
&\vv{CB}& = \ldots \ldots \vv{AC}
& \vv{AC}& = \ldots \ldots \vv{AB} 
\end{align*}  
\end{enumerate}
\end{exercise}

\begin{exercise} 	Placer les points $M$, $N$, $P$ pour chacun des cas suivants  :
\begin{enumerate}[label=\alph*),leftmargin=*]
\item $\vv{AM}=2\vv{AB}$; $\vv{AN}=-3\vv{AB}$ et $\vv{OP}=4\vv{AB}$

\vspace{0.5cm}

\begin{tikzpicture}[scale=0.8]
\tkzInit[xmax=6,ymax=1.5]%\tkzGrid[sub]
\tkzSetUpPoint[size=4, fill=black]
\tkzfctset{tan style/.style={->,>=latex, color=red, line width=1.2pt}}  

\def \xo{0.5}; 
\def \xa{4}; 
\def \xb{6};  
\tkzDefPoint(0,0){O}
\tkzDefShiftPoint[O](\xa,0){A} 
\tkzDefShiftPoint[O](\xb,0){B} 
\tkzDefPoint(12,0){C}  
\tkzDefPoint(-6,0){D}  
\tkzDrawSegments[-, >=latex, line width=1.2pt](C,D)


\tkzLabelPoints[above, font=\scriptsize](A,B,O) 

\foreach \i in {0,...,18}{%
%		\tkzMarkSegment[pos=\i/4,mark=|,  size=3pt](A,B)
%		\tkzMarkSegment[pos=\i/6,mark=|,  size=3pt](A,D)
\tkzMarkSegment[pos=\i/18,mark=|,  size=3pt](C,D)
} 
\end{tikzpicture}

\vspace{1cm}
\item $\vv{AM}=\tfrac{2}{5}\vv{AB}$; $\vv{BN}=-\tfrac{7}{5}\vv{AB}$ et $\vv{OP}=-\tfrac{1}{5}\vv{AB}$

\vspace{0.5cm}

\begin{tikzpicture}[scale=0.8]
\tkzInit[xmax=6,ymax=1.5]%\tkzGrid[sub]
\tkzSetUpPoint[size=4, fill=black]
\tkzfctset{tan style/.style={->,>=latex, color=red, line width=1.2pt}}  

\def \xo{2.5}; 
\def \xa{1}; 
\def \xb{6};  
\tkzDefPoint(0,0){O}
\tkzDefShiftPoint[O](\xa,0){A} 
\tkzDefShiftPoint[O](\xb,0){B} 
\tkzDefPoint(12,0){C}  
\tkzDefPoint(-6,0){D}  
\tkzDrawSegments[-, >=latex, line width=1.2pt](C,D)


\tkzLabelPoints[above, font=\scriptsize](A,B,O) 

\foreach \i in {0,...,18}{%
%		\tkzMarkSegment[pos=\i/4,mark=|,  size=3pt](A,B)
%		\tkzMarkSegment[pos=\i/6,mark=|,  size=3pt](A,D)
\tkzMarkSegment[pos=\i/18,mark=|,  size=3pt](C,D)
} 
\end{tikzpicture}	

\vspace{1cm}
\item $\vv{AM}=\tfrac{5}{3}\vv{AB}$; $\vv{BN}=-\tfrac{1}{6}\vv{AB}$ et $\vv{OP}=-\tfrac{4}{3}\vv{AB}$ 


\vspace{0.5cm}
\begin{tikzpicture}[scale=0.8]
\tkzInit[xmax=6,ymax=1.5]%\tkzGrid[sub]
\tkzSetUpPoint[size=4, fill=black]
\tkzfctset{tan style/.style={->,>=latex, color=red, line width=1.2pt}}  

\def \xo{5.0}; 
\def \xa{0}; 
\def \xb{6};  
\tkzDefPoint(\xo,0){O}
\tkzDefPoint(\xa,0){A} 
\tkzDefPoint(\xb,0){B} 
\tkzDefPoint(12,0){C}  
\tkzDefPoint(-6,0){D}  
\tkzDrawSegments[-, >=latex, line width=1.2pt](C,D)


\tkzLabelPoints[above, font=\scriptsize](A,B,O) 

\foreach \i in {0,...,18}{%
%		\tkzMarkSegment[pos=\i/4,mark=|,  size=3pt](A,B)
%		\tkzMarkSegment[pos=\i/6,mark=|,  size=3pt](A,D)
\tkzMarkSegment[pos=\i/18,mark=|,  size=3pt](C,D)
} 
\end{tikzpicture}	

\vspace{0.5cm}
\item $2\vv{AM}=-\vv{AB}$; $4\vv{BN}=3\vv{AB}$ et $2\vv{OP}=-3\vv{AB}$


\vspace{0.5cm}
\begin{tikzpicture}[scale=0.8]
\tkzInit[xmax=6,ymax=1.5]%\tkzGrid[sub]
\tkzSetUpPoint[size=4, fill=black]
\tkzfctset{tan style/.style={->,>=latex, color=red, line width=1.2pt}}  

\def \xo{0.0}; 
\def \xa{3}; 
\def \xb{7};  
\tkzDefPoint(\xo,0){O}
\tkzDefPoint(\xa,0){A} 
\tkzDefPoint(\xb,0){B}  
\tkzDefPoint(12,0){C}  
\tkzDefPoint(-6,0){D}  
\tkzDrawSegments[-, >=latex, line width=1.2pt](C,D)


\tkzLabelPoints[above, font=\scriptsize](A,B,O) 

\foreach \i in {0,...,18}{%
%		\tkzMarkSegment[pos=\i/4,mark=|,  size=3pt](A,B)
%		\tkzMarkSegment[pos=\i/6,mark=|,  size=3pt](A,D)
\tkzMarkSegment[pos=\i/18,mark=|,  size=3pt](C,D)
} 
\end{tikzpicture}	


\vspace{0.5cm}
\item $3\vv{AM}=2\vv{AB}$; $6\vv{BN}=11\vv{BA}$ et $6\vv{OP}=-\vv{BA}$


\vspace{0.5cm}
\begin{tikzpicture}[scale=0.8]
\tkzInit[xmax=6,ymin=-2,ymax=1.5]%\tkzGrid[sub]
\tkzSetUpPoint[size=4, fill=black]
\tkzfctset{tan style/.style={->,>=latex, color=red, line width=1.2pt}}  

\def \xo{5.0}; 
\def \xa{0}; 
\def \xb{6};  
\tkzDefPoint(\xo,0){O}
\tkzDefPoint(\xa,0){A} 
\tkzDefPoint(\xb,0){B}  
\tkzDefPoint(12,0){C}  
\tkzDefPoint(-6,0){D}  
\tkzDrawSegments[-, >=latex, line width=1.2pt](C,D)


\tkzLabelPoints[above, font=\scriptsize](A,B,O) 

\foreach \i in {0,...,18}{%
%		\tkzMarkSegment[pos=\i/4,mark=|,  size=3pt](A,B)
%		\tkzMarkSegment[pos=\i/6,mark=|,  size=3pt](A,D)
\tkzMarkSegment[pos=\i/18,mark=|,  size=3pt](C,D)
} 
\end{tikzpicture}	


\vspace{0.5cm}

%		\item $\vv{AP}=\frac{7}{2}\vv{AB}$
%		\item $\vv{AW}=-\frac{9}{2}\vv{AB}$
%		\item $\vv{AS}=-\frac{4}{5}\vv{AB}$
%		\item $\vv{AT}=-\frac{1}{4}\vv{AB}$
%		\item $\vv{AU}= \frac{12}{5}\vv{AB}$
%		\item $\vv{AV}= \frac{3}{2}\vv{AB}$
%		\item $\vv{BM}=\frac{12}{5}\vv{AB}$
%		\item $\vv{BN}=4\vv{AB}$
%		\item $\vv{BP}=\frac{8}{5}\vv{AB}$
%		\item $\vv{BQ}=-3\vv{AB}$
%		\item $\vv{BS}=-\frac{6}{5}\vv{AB}$
%		\item $\vv{BT}=-\frac{1}{4}\vv{AB}$ 
\end{enumerate}
\end{exercise} 
%\begin{tikzpicture}[scale=0.8]
%	\tkzInit[xmax=6,ymax=1.5]%\tkzGrid[sub]
%	\tkzSetUpPoint[size=4, fill=black]
%	\tkzfctset{tan style/.style={->,>=latex, color=red, line width=1.2pt}}  
%	
%	\def \xa{6}; 
%	\def \ya{0};  
%	\tkzDefPoint(0,0){A} \tkzDefShiftPoint[A](\xa,\ya){B} 
%	\tkzDefPoint(12,0){C}  
%	\tkzDefPoint(-6,0){D}  
%	\tkzDrawSegments[-, >=latex, line width=1.2pt](C,D)
%	
%	
%	\tkzLabelPoints[above, font=\scriptsize](A,B) 
%	
%	\foreach \i in {0,...,6}{%
%		\tkzMarkSegment[pos=\i/6,mark=|,  size=3pt](A,B)
%		\tkzMarkSegment[pos=\i/6,mark=|,  size=3pt](A,D)
%		\tkzMarkSegment[pos=\i/6,mark=|,  size=3pt](B,C)
%	} 
%\end{tikzpicture} 


\begin{example}\hfill   

\noindent	\begin{minipage}{1.0\columnwidth}\centering
\begin{tikzpicture}[scale=0.95]

\def \xmin{-7}; 
\def \xmax{11}; 

\def \ymin{-3};
\def \ymax{4};	
\def \alpha{75}

\pgftransformcm{1}{0}{1*cos(\alpha)}{1*sin(\alpha)}{\pgfpointorigin}
\path[clip,preaction = {draw=gray}] 
({\xmax-\ymin*cos(\alpha)} , {\ymin}) -- ({\xmax-\ymax*cos(\alpha)},{\ymax}) -- ({\xmin-\ymax*cos(\alpha)},{\ymax}) -- ({\xmin-\ymin*cos(\alpha)},{\ymin}) -- cycle; 
\tkzInit[xmin=-8, ymin=-8,  xmax=8.5,ymax=9.5]
\begin{scope}[dashed] 
\tkzGrid[ xstep=1, ystep=1](-11,-9)(13,15)
\end{scope}
%	\foreach \i in {-10, -9,...,10} {
%		\draw[ color =gray, dashed] (-11,-\i-10) -- (11,12-\i);
%		%		\tkzFct[domain = -10:10]{-\x-\i};
%	}

\def \xa{5}; 
\def \ya{0}; 

\def \xb{0};
\def \yb{2};

\def \xc{2};
\def \yc{3};

\tkzDefPoint(-7,3){U0} \tkzDefShiftPoint[U0](\xa,\ya){U1}
\tkzDefPoint(-3,-2){V0} \tkzDefShiftPoint[V0](\xb,\yb){V1} 
\tkzDefPoint(-6,-2){W0} \tkzDefShiftPoint[W0](\xc,\yc){W1}
\tkzDrawSegments[->,>=latex,line width =2 pt, color = Maroon](U0,U1 V0,V1  W0,W1)
\tkzLabelSegment[above](U0,U1){$\vv{v}$}
\tkzLabelSegment[left](V0,V1){$\vv{u}$}
\tkzLabelSegment[left](W1,W0){$\vv{w}$}

% 	\tkzDrawX[left=5pt, label=, font=\scriptsize, line width=1.2pt ] %, style=dashed
% 	\tkzDrawY[below =5pt, label=, font=\scriptsize, line width=1.2pt]  %, style=dashed


\tkzDefPoint(0,0){A}
\tkzDefPoint(4,-2){B}  
\tkzDrawPoints(A,B)
\tkzLabelPoints[font=\scriptsize](A,B) 

\end{tikzpicture} 
\end{minipage}  \\
Placer les points $A_0$, $A_1$,  $B_0$, $B_1$ tel que $\vv{AA_0}=\tfrac{3}{2}\vv{u}$; $\vv{AA_1}=-\vv{u}$; $\vv{BB_0}=\tfrac{3}{5}\vv{v}$ et $\vv{B_0B_1}=-\tfrac{7}{5}\vv{v}$.
%	\begin{enumerate}[label=\alph*),leftmargin=*]
%		\item Placer les points $A_0$, $A_1$ tel que $\vv{AA_0}=\tfrac{3}{2}\vv{u}$ et $\vv{AA_1}=-\vv{u}$. 
%		\item Placer les points $B_0$, $B_1$ tel que $\vv{BB_0}=\tfrac{3}{5}\vv{v}$ et $\vv{B_0B_1}=-\frac{7}{5}\vv{v}$.
%		\item Placer les points $C_0$, $C_1$, $C_2$ et $C_3$ tel que $\vv{AC_0}=\vv{w}$, $\vv{AC_1}=-\frac{1}{2}\vv{w}$,  $\vv{BC_2}=\frac{2}{3}\vv{w}$ et  $\vv{BC_3}=\frac{5}{3}\vv{w}$
%		\item  Compléter : $\vv{A_0A_1}=\ldots\ldots\vv{u}$; $\vv{B_0B_1}=\ldots\ldots\vv{v}$ et $\vv{C_2C_3}=\ldots\ldots\vv{w}$.
%	\end{enumerate} 
\end{example}
\begin{exercise} 
Pour chacun des cas suivants, dites s'il existe un réel $k$ tel que  $\vv{PQ}=k\vv{AB}$ puis préciser si les droites $(PQ)$ et $(AB)$ sont parallèles.
\begin{multicols}{2} 
\begin{enumerate}[label=\alph*),leftmargin=*]
\item $\vv{PQ}\begin{pmatrix}
4\\6
\end{pmatrix}$ et $\vv{AB}\begin{pmatrix}
-16\\-24
\end{pmatrix} $
\item $P(1\; ;\;-1)$, $Q(4\; ;\;3)$, $A(-1\; ;\;5)$ et $B(7\; ;\;1)$ 
\item $\vv{PQ}\begin{pmatrix}
-3\\5
\end{pmatrix}$ et $\vv{AB}\begin{pmatrix}
-9\\14
\end{pmatrix}$
%			\item $P(6;20)$, $Q(18;2)$ et $\vv{AB}\begin{pmatrix}
%				28\\-42
%			\end{pmatrix}$ 
\item $P(1\; ;\;4)$, $Q(-3\; ;\;5)$, $A(5\; ;\;7)$ et $B(9\; ;\;6)$ 
%			\item $P(1\; ;\;-1)$, $Q(-2\; ;\; 3)$, $A(3\; ;\;1)$ et $B(-3\; ;\;9)$ 
%			\item $P(-1\; ;\;5)$, $Q(3\; ;\; -4)$, $A(1 \; ;\;1)$ et $B(9\; ;\;-17)$ 
%			\item $P(2\; ;\;3)$, $Q(-1\; ;\; 3)$, $A(-2\; ;\;3)$ et $B(-11\; ;\;3)$ 
%			\item $P(3\; ;\;-4)$, $Q(3\; ;\; 5)$, $A(3\; ;\;-1)$ et $B(3\; ;\;-28)$  
%			\item $P(−2;1)$, $Q(1;3)$, $A(4;5)$ et $B(13;-13)$
\end{enumerate}
\end{multicols}
\end{exercise}
\begin{example}[Je fais] Dans un repère $(O;~I,J)$ on définit $A(-15~;~12)$ et $B(8~;~-4)$.\\ Sachant que $\vv{CO}=10\vv{AB}$, déterminer les coordonnées de $C$.

\vfill
\end{example}
\begin{exercise} Mêmes consignes
\begin{enumerate}[label=\alph*),leftmargin=*]
\item $A(15~;~-10)$ et $B(-1~;~-4)$ et $\vv{OC}=-11\vv{AB}$. %176;-66
\item $A(1~;~-16)$ et $B(-9~;~9)$ et $\vv{CB}=-14\vv{AB}$.%-149;359
\item  $A(5~;~9)$ et $B(-13~;~-14)$ et $\vv{CB}=-5\vv{AB}$.%-103;-129
\end{enumerate} 
\end{exercise}

%----------------------------------------------------------------------------------------
%	 
%----------------------------------------------------------------------------------------

\newpage 
\loadgeometry{marginlayout}  
\pagestyle{scrheadings} % Print headers again  
\KOMAoptions{
headwidth=textwithmarginpar,
headsepline=true,
headsepline= 0.5pt,
footwidth=textwithmarginpar,
} 
\onehalfspacing

\section{Colinéarité de deux vecteurs}
\begin{definition}$\vv{u}\neq \vv{0}$.  Les vecteurs colinéaires à $\vv{u}$ sont tous les multiples $k\vv{u}$ ou $k\in\mathbb{R}$.\\
Il s'agit du vecteur nul $\vv{0}$ et de tous les vecteurs de même direction que $\vv{u}$. 
\end{definition} 
\begin{tBox} Deux vecteurs $\vv{u}$ et $\vv{v}$  sont \textbf{colinéaires}\fg{} (en abrégé $\vv{u}\propto\vv{v}$) si l'un est un multiple de l'autre. 
\end{tBox} 
%Les vecteurs colinéaires n'ont pas forcément la même longueur, ni le même sens.
%\begin{tBox} Deux vecteurs $\vv{u}$ et $\vv{v}$ \textbf{ne sont pas colinéaires} signifie :
%\begin{enumerate}[label=\alph*)]
%\item $\vv{u}\neq \vv{0}$ et $\vv{v}\neq \vv{0}$
%\item ils ont des directions différentes
%\end{enumerate}
%\end{tBox}

\begin{postulatT} trois points $A$, $B$ et $C$ sont alignés si et seulement si deux parmi les vecteurs $\vv{AB}$, $\vv{AC}$ ou $\vv{BC}$ sont colinéaires.
\end{postulatT}
\begin{postulatT} les droites $(AB)$ et  $(CD)$ sont parallèles si et seulement si  les vecteurs $\vv{AB}$ et $\vv{CD}$ sont colinéaires.
\end{postulatT}
\begin{definition}
Dans un repère quelconque $(O;\vec{\imath},\vec{\jmath})$. Soit deux vecteurs $\vv{u}\begin{pmatrix}
x\\y
\end{pmatrix}$ et $\vv{v}\begin{pmatrix}
x'\\y'
\end{pmatrix}$.\\
On appelle \og déterminant de $\vv{u}$ et $\vv{v}$ \fg{} le nombre noté :
\setlength{\abovedisplayskip}{3pt}
\setlength{\belowdisplayskip}{3pt} 
\[
\det(\vv{u}\; ; \; \vv{v}) = \mdet{x & x' \\ y & y'} = xy'- x'y
\]
\end{definition}
\begin{theorem}[Colinéarité à l'aide des coordonnées] Deux vecteurs $\vv{u}$ et $\vv{v}$ sont colinéaires si et seulement si $\det(\vv{u};\vv{v})=0$.
\end{theorem}
\begin{proof}\hfill \vfill 
\end{proof}


%----------------------------------------------------------------------------------------
%	EXERCICES
%----------------------------------------------------------------------------------------

\newpage 
\loadgeometry{widelayout}  
\pagestyle{scrheadings} % Print headers again  
\KOMAoptions{
headwidth=textwithmarginpar,
headsepline=true,
headsepline= 0.5pt,
footwidth=textwithmarginpar,
} 
\onehalfspacing
\subsection*{Exercices : Colinéarité de vecteurs}
\begin{example}
Calculer les déteriminant des vecteurs suivants :
\begin{enumerate}[label=\alph*),leftmargin=*]
\item \rule{0pt}{2em}
$\vv{u}\begin{pmatrix}
3\\2
\end{pmatrix}$ et $\vv{v}\begin{pmatrix}
4\\5
\end{pmatrix}$. \quad $\det(\vv{u};\vv{v})=\mdet{3 & 4 \\ 2 & 5}=$\vfill $\det(\vv{v};\vv{u})=$
\item  \rule{0pt}{2em}
$\vv{u}\begin{pmatrix}
-3\\2
\end{pmatrix}$ et $\vv{v}\begin{pmatrix}
-1\\4
\end{pmatrix}$.\quad $\det(\vv{u};\vv{v})=$
\item  \rule{0pt}{2em}
$\vv{u}\begin{pmatrix}
-4\\5
\end{pmatrix}$ et $\vv{v}\begin{pmatrix}
8\\-10
\end{pmatrix}$. 
\end{enumerate}
\end{example}

\begin{example}[Interprétation géométrique du déterminant dans un repère orthonormé]\hfill 

Soit un repère \textbf{orthonormé} $(O;I,J)$, et les points $M(a,b)$ et $N(c,d)$ et $P$ tel que $OPMN$ est un parallélogramme. Pour simplifier on suppose que $a$, $b$, $c$ et $d$  sont des réels strictement positifs.

\noindent\parbox{0.5\linewidth}{
\begin{tikzpicture}[scale=0.8]
\tkzSetUpPoint[size=3, fill=black]
\tkzfctset{tan style/.style={->,>=latex, color=red, line width=1.2pt}}   
\tkzSetUpLine[line width= 1.2pt, add=.2 and .2] 
\tkzInit[xmin=-0.5,xmax=7.5,ymin=-0.5,ymax=6.5  ]
\tkzDrawX[  noticks , label=$x$,  font=\small, line width=1.2pt ] %, style=dashed
\tkzDrawY[  noticks , label=$y$,  font=\small, line width=1.2pt]  %, style=dashed
\tkzRep[color  = Maroon,line width = 2pt, posxlabel=below, posylabel=left, xnorm=1,ynorm=1, xlabel={}, ylabel={}]

\tkzDefPoint(0,0){O};
\tkzDefPoint(2,4){N};
\tkzDefPoint(2,6){xN};
\tkzDefPoint(0,4){yN};
\tkzDefPoint(5,2){M};
\tkzDefPoint(5,0){xM};
\tkzDefPoint(7,2){yM};
\tkzDefPoint(7,6){P};
\tkzDefPoint(7,0){xP};
\tkzDefPoint(0,6){yP};
\tkzDrawPolygon[fill=gray!20](O,M,P,N);
\tkzDrawSegments[dashed](N,xN);
\tkzDrawSegments[dashed](M,yM);
%	\tkzDrawSegments[dashed](P,xP P,yP);
\tkzPointShowCoord(P)
\tkzPointShowCoord[noydraw](M)
\tkzPointShowCoord[noxdraw](N)
\tkzDrawSegment[thick,->, >=latex](O,M);
\tkzDrawSegment[->,  >=latex  ](O,N);
\tkzDrawSegment[- , dashed, >=latex](M,P);
\tkzLabelPoint[above left,font=\small](M){$M(a,b)$}
\tkzLabelPoint[below right,font=\small](N){$N(c,d)$}
\tkzLabelPoint[above,font=\small](P){$P(\qquad,\qquad)$}
\tkzLabelPoints[below left,font=\small](O)
\end{tikzpicture} 
}

\end{example}
\begin{example} Déterminer si les vecteurs $\vv{u}$ et $\vv{v}$ sont colinéaires. 
\begin{multicols}{2}
\begin{enumerate}[label=\alph*),leftmargin=*]
\item \rule{0pt}{2em}
$\vv{u}\begin{pmatrix}
2\\3
\end{pmatrix}$ et $\vv{v}\begin{pmatrix}
6\\9
\end{pmatrix}$. 
\item \rule{0pt}{2em}
$\vv{u}\begin{pmatrix}
\sqrt{3}+\sqrt{2}\\1
\end{pmatrix}$ et $\vv{v}\begin{pmatrix}
1\\\sqrt{3}-\sqrt{2}
\end{pmatrix}$. 
\end{enumerate}
\end{multicols}
\end{example}
\vfill 
\begin{example}
Soient un repère orthonormé et les points $A(0~;~-2)$, $B(1~;~1)$ et $C(2~;~-1)$. Calculer l'aire du triangle $ABC$.
\end{example}
\vfill 
\begin{example}
Montrer que les points $A(2~;~5)$, $B(3~;~8)$ et $C(-5~;~-16)$ sont alignés.
\end{example}
\vfill 

\newpage
\begin{example}
Soient $A(1~;~3)$, $B(5~;~-2)$, $C(-1~;~6)$ et $D(7~;~-4)$. 
Montrer que les droites (AB) et (CD) sont parallèles.
\end{example}
\vfill 


\begin{multicols}{2} 
\begin{exercise}


Dans chaque cas, dire si les vecteurs sont colinéaires. 
%	\begin{multicols}{3}
\begin{enumerate}[label=\alph*),leftmargin=*]
\item  $\vv{u}\begin{pmatrix}
2\\-3 
\end{pmatrix}$ et $\vv{v}\begin{pmatrix}
-1\\\tfrac{3}{2} 
\end{pmatrix}$  
\vspace{0.5cm}
\item  $\vv{u}\begin{pmatrix}
\tfrac{1}{2}\\\tfrac{1}{3} 
\end{pmatrix}$ et $\vv{v}\begin{pmatrix}
\tfrac{4}{5}\\ -\tfrac{1}{3}
\end{pmatrix}$  
\vspace{0.5cm}
\item  $\vv{u}\begin{pmatrix}
\sqrt{2}\\\sqrt{3} 
\end{pmatrix}$ et $\vv{v}\begin{pmatrix}
2\\ -\sqrt{6}
\end{pmatrix}$  
\vspace{0.5cm} 
\end{enumerate}
%\end{multicols}
\end{exercise}
\begin{exercise}Dans chaque cas les vecteurs $\vv{u}$ et $\vv{v}$ sont colinéaires. Écrire une équation vérifiée par $m$ et la résoudre.   
%		\begin{multicols}{3}
\begin{enumerate}[label=\alph*),leftmargin=*]
\item  $\vv{u}\begin{pmatrix}
2\\6
\end{pmatrix}$ et $\vv{v}\begin{pmatrix}
m\\ 3
\end{pmatrix}$  
\vspace{0.5cm}
\item  $\vv{u}\begin{pmatrix}
-m\\4m-3
\end{pmatrix}$ et $\vv{v}\begin{pmatrix}
1\\ -3
\end{pmatrix}$  
\vspace{0.5cm}
\item  $\vv{u}\begin{pmatrix}
27\\2m
\end{pmatrix}$ et $\vv{v}\begin{pmatrix}
2m\\ 3
\end{pmatrix}$  
\vspace{0.5cm} 
\end{enumerate} 
%\end{multicols}
\end{exercise}
\begin{exercise} 
Dans un repère orthonormé $( O ; \vec{\imath} , \vec{\jmath} )$.
\begin{enumerate}[label=\alph*),leftmargin=*]
\item Montrer que $A\left(-2;1\right)$, $B\left(3;3\right)$ , $C\left(1;\tfrac{11}{5}\right)$ sont alignés.
%		\vfill
\item Trouver $a$ pour que $A(1 ; 2)$, $B(-3 ; 0)$ et $C(-14 ; a)$ soient alignés.
%		\vfill 
\end{enumerate} 
\end{exercise}  
\begin{exercise} 
Dans un repère, soit les points: $A(-2~;~1)$, $B(3~;~3)$, $C\left(1~;~\tfrac{11}{5}\right)$ et $D\left(\tfrac{45}{2}~;~\tfrac{54}{5}\right)$
\begin{enumerate}[label=\alph*),leftmargin=*]
\item Démontrer que les points $A$, $B$ et $C$ sont alignés. 
%		\vspace{2cm}
\item Les points $A$, $B$ et $D$ sont-ils alignés ?
%		\vspace{2cm}
\end{enumerate}

\end{exercise}
%\newpage 

\begin{exercise} 
Pour chaque cas, justifier si les points donnés sont alignés.
\begin{enumerate}[label=\alph*)]
\item $A(5;10)$, $B(-5;5)$ et $C(-1;7)$
\item $D(-3;-4)$, $E(-6;-6)$ et $F(1;-1)$
\item $K(-9;4)$, $L(0;3)$ et $M(9,2)$
\item $P(-3;2)$, $Q(-6;-2)$ et $R(-10,-6)$
\end{enumerate}  
\end{exercise}
\begin{exercise}  
Dans un repère, on donne les points : $M(0~;~-3)$, $N(2~;~3)$, $P(-9~;~0)$ et $Q(-1~;~-1)$
\begin{enumerate}[label=\arabic*),leftmargin=*]
\item Calculer les coordonnées des points $A$ et $B$ tels que: 
\begin{enumerate}[label=\alph*),leftmargin=*]
\item $\vv{NA}=\dfrac{1}{2}~\vv{MN}$
%			\vspace{2cm}
\item $\vv{MB}=3~\vv{MQ}$
%			\vspace{2cm}
\end{enumerate}
\item Calculer les coordonnées des vecteurs $\vv{PA}$ et $\vv{PB}$.
%		\vspace{2cm} 
\item Démontrer que les points $P$, $A$ et $B$ sont alignés.
%		\vspace{3cm}
\end{enumerate}
\end{exercise}
\begin{exercise} 
Soit les points  $A(1 ; 2)$, $B(-3 ; 0)$ et $C(-7 ; a)$ dans un repère orthonormé. Trouver 2 valeurs de $a$ pour lesquelles le triangle $ABC$ est d'aire $12$.  
\end{exercise}
\end{multicols}
%\vfill 
%\begin{exercise} 
%	Dans un repère, on donne les points: $A(1~;~-1)$, $B(-1~;~-2)$ et $C(-2~;~2)$
%	\begin{enumerate}[label=\arabic*),leftmargin=*]
%		\item Déterminer les coordonnées du point $G$ vérifiant $\vv{GA}+2\vv{GB}+\vv{GC}=\vv{0}$.
%%		\vspace{2.5cm} 
%		\item Déterminer les coordonnées du point $D$ vérifiant $\vv{BD}=\vv{BA}+\vv{BC}$.
%%		\vspace{2.5cm} 
%		\item Faire une figure. Que peut-on conjecturer pour les points $B$, $G$ et $D$ ?  
%		Démontrer cette conjecture.
%	\end{enumerate} 
%\end{exercise}

%\vfill 





%----------------------------------------------------------------------------------------
%	 
%----------------------------------------------------------------------------------------

\newpage 
\loadgeometry{marginlayout}  
\pagestyle{scrheadings} % Print headers again  
\KOMAoptions{
	headwidth=textwithmarginpar,
	headsepline=true,
	headsepline= 0.5pt,
	footwidth=textwithmarginpar,
} 
\onehalfspacing


\section{Somme de vecteurs}
\begin{theorem}
L'enchainement d'une translation de vecteur $\vv{u}$ puis d'une translation de vecteur $\vv{v}$ est aussi une translation.

Le vecteur $\vv{u}+\vv{v}$ est le vecteur associé à la translation obtenue.
\end{theorem} 
%\begin{remark} C'est la conséquence du théorème  \ref{thm.transitif.para} 
%\end{remark}  

\begin{figure}[!h]
\begin{tikzpicture}[scale=1.3]
\tkzSetUpPoint[size=3, fill=black]
\tkzfctset{tan style/.style={->,>=latex, color=red, line width=1.2pt}}   
\tkzSetUpLine[line width= 1.2pt, add=.2 and .2] 

\tkzInit[xmax=6,ymax=1.5]%\tkzGrid[sub]   

\def \xa{0.8}; 
\def \ya{2.4}; 

\def \xb{4};
\def \yb{0};

\tkzDefPoint(-0.7,0){U0} \tkzDefShiftPoint[U0](\xa,\ya){U1}
\tkzDefPoint(1.0,3.5){V0} \tkzDefShiftPoint[V0](\xb,\yb){V1}

\tkzDrawSegments[->, >=latex, line width=1.2pt](U0,U1 V0,V1)
\tkzLabelSegment(U0,U1){\small $\vv{u}$}
\tkzLabelSegment(V0,V1){\small $\vv{v}$}

\tkzDefPoint(0.8,0){A}  
\tkzDefPointBy[translation=from U0 to U1](A)\tkzGetPoint{B}
\tkzDefPointBy[translation=from V0 to V1](B)\tkzGetPoint{C}

\tkzDrawSegments[->, >=latex, line width=0.6pt, dashed](A,B)
\tkzDrawSegments[-, line width=0.6pt, dashed](B,C)
\tkzDrawSegments[->, >=latex, line width=1.2pt](A,C)

\tkzLabelPoint[below](A){\scriptsize$A$}
\tkzLabelPoint[above left](B){\scriptsize$B$}
\tkzLabelPoint[right](C){\scriptsize$C$}
\tkzLabelSegment[left](A,B){\small$\vv{AB}$}
\tkzLabelSegment[above](B,C){\small$\vv{BC}$}
\tkzLabelSegment[below, sloped](A,C){\small$\vv{AC}=\vv{AB}+\vv{BC}$}

\tkzDefPoint(2.0,-3.5){A0}  
\tkzDefPointBy[translation=from U0 to U1](A0)\tkzGetPoint{B0}
\tkzDefPointBy[translation=from V0 to V1](B0)\tkzGetPoint{C0}

\tkzDrawSegments[->, >=latex, line width=0.6pt, dashed](A0,B0)
\tkzDrawSegments[-, line width=0.6pt, dashed](B0,C0)
\tkzDrawSegments[->, >=latex, line width=1.2pt](A0,C0)

\tkzLabelPoint[below](A0){\scriptsize$X$}
\tkzLabelPoint[above left](B0){\scriptsize$Y$}
\tkzLabelPoint[right](C0){\scriptsize$Z$}
\tkzLabelSegment[left](A0,B0){\small$\vv{XY}$}
\tkzLabelSegment[above](B0,C0){\small$\vv{YZ}$}
\tkzLabelSegment[below right](A0,C0){\small$\vv{XZ}$} 
\end{tikzpicture}
\caption{ 
$C$ est l'image de $A$ par l'enchaînement des translation de vecteur $\protect\vv{u}$ puis $\protect\vv{v}$. \\
$\protect\vv{u}+\protect\vv{v} = \protect\vv{AB}+\protect\vv{BC}=\protect\vv{AC}$.\\
En prenant un autre point de départ $A'$ on obtient un autre représentant de $\protect\vv{u}+\protect\vv{v}$ :
$\protect\vv{XZ}= \protect\vv{AC}=\protect\vv{u}+\protect\vv{v}$
}
\end{figure}

\begin{tBox}[frametitle={Relation de Chasles}]
Pour tout points $A$, $B$ et $C$ :
\(  \vv{AC}=\vv{AB}+\vv{BC} \) 
\end{tBox}   
\begin{example}Placer les points $B$ et $Q$ tel que $\protect\vv{AB}=\protect\vv{u}+\protect\vv{v}$ et $\protect\vv{PQ}=\protect\vv{v}+\protect\vv{w}$


\end{example} 
\begin{figure}[!h]\centering 
\begin{tikzpicture}[scale=0.8] 

\tkzSetUpPoint[size=3, fill=black]
\tkzfctset{tan style/.style={->,>=latex, color=red, line width=1.2pt}}   
\tkzSetUpLine[line width= 1.2pt, add=.2 and .2] 

\def \xmin{-6}; 
\def \xmax{7.5}; 
\def \ymin{-3};
\def \ymax{5};	
\def \alpha{-00}
\def \beta{80}
\def \xnorm{1.2};
\def \ynorm{0.8}; 
\pgftransformcm{\xnorm*cos(\alpha)}{\xnorm*sin(\alpha)}{\ynorm*cos(\beta)}{\ynorm*sin(\beta)}{\pgfpointorigin}
\path[clip,preaction = {draw=gray}] 
({(\xmax*\ynorm*sin(\beta)-\ymax*\ynorm*cos(\beta))/(\xnorm*\ynorm*sin(\beta-\alpha))},{(-\xmax*\xnorm*sin(\alpha)+\ymax*\xnorm*cos(\alpha))/(\xnorm*\ynorm*sin(\beta-\alpha))}	) -- ( 
{(\xmax*\ynorm*sin(\beta)-\ymin*\ynorm*cos(\beta))/(\xnorm*\ynorm*sin(\beta-\alpha))},{(-\xmax*\xnorm*sin(\alpha)+\ymin*\xnorm*cos(\alpha))/(\xnorm*\ynorm*sin(\beta-\alpha))}) -- (
{(\xmin*\ynorm*sin(\beta)-\ymin*\ynorm*cos(\beta))/(\xnorm*\ynorm*sin(\beta-\alpha))},{(-\xmin*\xnorm*sin(\alpha)+\ymin*\xnorm*cos(\alpha))/(\xnorm*\ynorm*sin(\beta-\alpha))}) -- (
{(\xmin*\ynorm*sin(\beta)-\ymax*\ynorm*cos(\beta))/(\xnorm*\ynorm*sin(\beta-\alpha))},{(-\xmin*\xnorm*sin(\alpha)+\ymax*\xnorm*cos(\alpha))/(\xnorm*\ynorm*sin(\beta-\alpha))}) --   cycle; 
\tkzInit[xmin=-8, ymin=-8,  xmax=8.5,ymax=9.5]
\begin{scope}[dashed] 
\tkzGrid[ xstep=1, ystep=1](-11,-9)(10,15)
\end{scope}  
%		\tkzDrawX[left=5pt, label=, font=\scriptsize, line width=1.2pt ]  
%		\tkzDrawY[below =5pt, label=, font=\scriptsize, line width=1.2pt] 
%		\tkzRep[color  = Maroon,xlabel=,ylabel=, line width = 2pt,xnorm=1, ynorm=1]  
%		\tkzDefPoint(0,0){O} 
%		\tkzDefPoint(1,0){I} 
%		\tkzDefPoint(0,1){J} 
%		\tkzLabelPoints[below left,  font=\scriptsize](O)
%		\tkzLabelPoints[below,  font=\scriptsize](I)
%		\tkzLabelPoints[left,  font=\scriptsize](J)

\def \xa{2}; 
\def \ya{3}; 

\def \xb{4};
\def \yb{1};

\def \xc{-3};
\def \yc{0};

\tkzDefPoint(-4,-2){U0} \tkzDefShiftPoint[U0](\xa,\ya){U1}
\tkzDefPoint(-5,3){V0} \tkzDefShiftPoint[V0](\xb,\yb){V1}
\tkzDefPoint(5,5){W0} \tkzDefShiftPoint[W0](\xc,\yc){W1}
\tkzDrawSegments[->,>=latex,line width =2 pt, color = Maroon](U0,U1 V0,V1 W0,W1)
\tkzLabelPoint[below, font=\small](U0){$A$}
\tkzDrawPoint(4,-2)
\tkzDrawPoint(U0)
\tkzLabelPoint[below, font=\small](4,-2){$P$}
\tkzLabelSegment[font=\small,above](U0,U1){$\vv{v}$}
\tkzLabelSegment[font=\small,above](V0,V1){$\vv{u}$}
\tkzLabelSegment[font=\small,above](W0,W1){$\vv{w}$}
\end{tikzpicture}
%		\caption{Placer les points $B$ et $Q$ tel que $\protect\vv{AB}=\protect\vv{u}+\protect\vv{v}$ et $\protect\vv{PQ}=\protect\vv{v}+\protect\vv{w}$}
\end{figure} 
 
\begin{definition}[Opposé du vecteur]
se note $-\vv{u}$ et vérifie :
\setlength{\abovedisplayskip}{3pt}
\setlength{\belowdisplayskip}{3pt} 
$$\vv{u}+(-\vv{u}) =\vv{0}$$ 
Deux vecteurs opposés ont même direction, même longueur, mais sont de sens contraires.
\end{definition}
\begin{example}Pour les points $A$ et $B$, $\vv{BA}$ est l'opposé du vecteur $\vv{AB}$ :
\setlength{\abovedisplayskip}{3pt}
\setlength{\belowdisplayskip}{3pt} 
$$\vv{AB}+\vv{BA}=\vv{AA}=\vv{0}$$  
On écrira : $\vv{BA}=-\vv{AB}$
\end{example}
\begin{tBox}
Soit un repère $(O;I,J)$ et les vecteurs $\vv{v}\begin{pmatrix}
x\\y
\end{pmatrix}$ et $\vv{u}\begin{pmatrix}
x'\\y'
\end{pmatrix}$.\\
Le vecteur $\vv{u}+\vv{v}$ a pour coordonnées 
$(\vv{u}+\vv{v})\begin{pmatrix}
x+x'\\y+y '
\end{pmatrix}$\\
Le vecteur $-\vv{u}\begin{pmatrix}
-x\\-y
\end{pmatrix}$
\end{tBox}
\begin{figure}[!h]
	\begin{tikzpicture}[scale=0.85]
		\tkzSetUpPoint[size=3, fill=black]
		\tkzfctset{tan style/.style={->,>=latex, color=red, line width=1.2pt}}   
		\tkzSetUpLine[line width= 1.2pt, add=.5 and .5] 
		
		\tkzInit[xmin=-1, ymin=-1,  xmax=7,ymax=7]  
		\begin{scope}[dashed] 
			\tkzGrid 
		\end{scope}  
		\tkzDrawX[   label=$x$, font=\scriptsize, line width=1.2pt ] %, style=dashed
		\tkzDrawY[  label=$y$, font=\scriptsize, line width=1.2pt]  %, style=dashed
		\tkzRep[color  = Maroon,line width = 2pt, posxlabel=below, posylabel=left, xnorm=1,ynorm=1, xlabel={$\overrightarrow{\imath}$}, ylabel={$\overrightarrow{\jmath}$}]
		\tkzDefPoint(0,0){O}  
		\tkzDefPoint(1,0){I}  
		\tkzDefPoint(0,1){J}
		\tkzDefPoint(2,6){P}
		\tkzLabelPoint[below left](O){$O$}
		\tkzLabelPoint[above right](P){$P$}
		\tkzDefShiftPoint[O](2,1){U0}
		\tkzDefShiftPoint[U0](4,3){U1}
		\tkzDrawSegment[->,>=latex,line width=2pt,color=red](U0,U1)
		\tkzDrawSegment[->,>=latex,line width=1.2pt,color=blue](O,P)
%		\tkzLabelAngle[pos=1.25,font=\scriptsize](U,O,V){$\tfrac{\pi}{4}$}
%		\tkzMarkAngle[size=0.6](U,O,V)
		\tkzLabelSegment[above,sloped](U0,U1){$\vv{u}$}
		\tkzLabelSegment[above,sloped](O,P){$\vv{OP}$}
		\tkzLabelX[orig=false, 	font=\scriptsize]  
		\tkzLabelY[orig=false,  font=\scriptsize] 
	\end{tikzpicture} 
\caption{Décomposition de vecteurs dans la base $(\protect\vv{\imath} ~;~\protect\vv{\jmath})$ \protect\newline $\protect\vv{OP}=2 \protect\vv{\imath}+ 6\protect\vv{\jmath} $\protect\newline 
$\protect\vv{u}=4 \protect\vv{\imath}+ 3\protect\vv{\jmath} $
}
\end{figure}
\begin{definition} Soit trois points non alignés $O$, $I$ et $J$. Les vecteurs $\vv{OI}=\vv{\imath}$ et $\vv{OJ}=\vv{\jmath}$ ne sont pas colinéaires et forment une \textbf{base}.\\
	Si on muni le plan du repère   $(O~;~I~,~J)$ (noté encore $(O~;~\vv{OI}~,~\vv{OJ})$ ou $(O~;~\vv{\imath}~,~\vv{\jmath})$) alors :
\begin{enumerate}[leftmargin=*,label=\textbullet]
	\item Tout point $P$ admet un unique couple de coordonnées $(x_P~;~y_P)$   :
	\setlength{\abovedisplayskip}{3pt}
	\setlength{\belowdisplayskip}{3pt}
	$$\vv{OP}=x_P\vv{\imath}+y_P\vv{\jmath}$$  
	\item \noindent Tout vecteur $\vv{u}$ admet un unique couple de coordonnées $\begin{pmatrix}
		x\\ y
	\end{pmatrix}$ dans la \textbf{base} $(\vv{\imath}~;~\vv{\jmath})$ tel que 
	$\vv{u}=x\vv{\imath}+y\vv{\jmath}$ 
\end{enumerate}  
\end{definition} 
%----------------------------------------------------------------------------------------
%	EXERCICES
%----------------------------------------------------------------------------------------

\newpage 
\loadgeometry{widelayout}  
\pagestyle{scrheadings} % Print headers again  
\KOMAoptions{
headwidth=textwithmarginpar,
headsepline=true,
headsepline= 0.5pt,
footwidth=textwithmarginpar,
} 
\onehalfspacing
\section{Exercices : sommes de vecteurs}


\begin{example} Placer les points $B$, $D$ et $F$ tel que $\protect\vv{AB}=\protect\vv{u}+\protect\vv{v}$, $\protect\vv{CD}=\protect\vv{u}-\protect\vv{v}$ et $\protect\vv{EF}=\protect\vv{v}-\protect\vv{u}$ 


\end{example} 
\begin{figure}[h]\centering 
\begin{tikzpicture}[scale=0.8]

\def \xmin{-7}; 
\def \xmax{11}; 

\def \ymin{-3};
\def \ymax{4};	
\def \alpha{75}

\pgftransformcm{1}{0}{1*cos(\alpha)}{1*sin(\alpha)}{\pgfpointorigin}
\path[clip,preaction = {draw=gray}] 
({\xmax-\ymin*cos(\alpha)} , {\ymin}) -- ({\xmax-\ymax*cos(\alpha)},{\ymax}) -- ({\xmin-\ymax*cos(\alpha)},{\ymax}) -- ({\xmin-\ymin*cos(\alpha)},{\ymin}) -- cycle; 
\tkzInit[xmin=-8, ymin=-8,  xmax=17.5,ymax=9.5]
\begin{scope}[dashed] 
\tkzGrid[ xstep=1, ystep=1](-11,-9)(18,15)
\end{scope}
%	\foreach \i in {-10, -9,...,10} {
%		\draw[ color =gray, dashed] (-11,-\i-10) -- (11,12-\i);
%		%		\tkzFct[domain = -10:10]{-\x-\i};
%	}

\def \xa{4}; 
\def \ya{0}; 

\def \xb{0};
\def \yb{3};

\tkzDefPoint(-5,3){U0} \tkzDefShiftPoint[U0](\xa,\ya){U1}
\tkzDefPoint(-6,-1){V0} \tkzDefShiftPoint[V0](\xb,\yb){V1}
\tkzDrawSegments[->,>=latex,line width =2 pt, color = Maroon](U0,U1 V0,V1)
\tkzLabelSegment[above](U0,U1){$\vv{u}$}
\tkzLabelSegment[left](V0,V1){$\vv{v}$}

% 	\tkzDrawX[left=5pt, label=, font=\scriptsize, line width=1.2pt ] %, style=dashed
% 	\tkzDrawY[below =5pt, label=, font=\scriptsize, line width=1.2pt]  %, style=dashed


\tkzDefPoint(-1,-2){A}
\tkzDefPoint(5,2){C} 
\tkzDefPoint(10,-1){E}
\tkzDrawPoints(A,C,E)
\tkzLabelPoints[font=\scriptsize](A,C,E) 

\end{tikzpicture} 
%	\caption{Placer les points $B$, $D$ et $F$ tel que $\protect\vv{AB}=\protect\vv{u}+\protect\vv{v}$, $\protect\vv{CD}=\protect\vv{u}-\protect\vv{v}$ et $\protect\vv{EF}=\protect\vv{v}-\protect\vv{u}$}
\end{figure}


\noindent Pour les exercices 1 à 4,  $ABCDEF$ est un hexagone régulier de centre $O$ et rayon $1$. Écrire le plus simplement possible les sommes ci-dessous en utilisant des points de la figure ou le vecteur nul, ou un nombre.
\begin{exercise}\hfill 

\noindent
\parbox{0.40\linewidth}{\centering
\begin{tikzpicture} 
\tkzSetUpPoint[size=3, fill=black]
\tkzfctset{tan style/.style={->,>=latex, color=red, line width=1.2pt}}   
\tkzSetUpLine[line width= 1.2pt, add=.2 and .2] 

\tkzDefPoint(0,0){O}
\def \rayon{3.0};  
\tkzDefShiftPoint[O](0:\rayon){A}
\tkzDefShiftPoint[O](60:\rayon){B}
\tkzDefShiftPoint[O](120:\rayon){C}
\tkzDefShiftPoint[O](180:\rayon){D}
\tkzDefShiftPoint[O](240:\rayon){E}
\tkzDefShiftPoint[O](300:\rayon){F}
\tkzDrawPolygon(A,B,C,D,E,F)
\tkzLabelPoint[right](A){$A$}
\tkzLabelPoint[above right](B){$B$}
\tkzLabelPoint[above left](C){$C$}
\tkzLabelPoint[left](D){$D$}
\tkzLabelPoint[below left](E){$E$}
\tkzLabelPoint[below right](F){$F$}
\tkzDrawSegments[dashed, line width=1pt](O,A O,B O,C O,D O,E O,F) 
\end{tikzpicture} 
}\hfill 
\parbox{0.55\linewidth}{
\begin{enumerate}[label=\alph*)]
\item\rule{0pt}{2.1em} $\vv{OA}+\vv{AB}=$
\item\rule{0pt}{2.1em} $\vv{OA}+\vv{OC}=$
\item\rule{0pt}{2.1em} $\vv{OA}+\vv{OB}=$
\item\rule{0pt}{2.1em} $OA+OV=$
\item\rule{0pt}{2.1em} $\vv{EF}+\vv{FA}+\vv{AB}=$
\item\rule{0pt}{2.1em} $\vv{CD}+\vv{FA}=$
%	\item\rule{0pt}{2.1em} $\vv{OB}-\vv{OA}=$
%	\item\rule{0pt}{2.1em} $OB-OA=$
%	\item\rule{0pt}{2.1em} $\vv{CB}-\vv{FA}=$
%	\item\rule{0pt}{2.1em} $\vv{DO}-\vv{AB}=$
\end{enumerate}
}
\end{exercise}
\vfill 
\begin{exercise}\hfill  

\noindent
\parbox{0.40\linewidth}{\centering
\begin{tikzpicture} 
\tkzSetUpPoint[size=3, fill=black]
\tkzfctset{tan style/.style={->,>=latex, color=red, line width=1.2pt}}   
\tkzSetUpLine[line width= 1.2pt, add=.2 and .2] 

\tkzDefPoint(0,0){O}
\def \rayon{3.0};  
\tkzDefShiftPoint[O](0:\rayon){A}
\tkzDefShiftPoint[O](60:\rayon){B}
\tkzDefShiftPoint[O](120:\rayon){C}
\tkzDefShiftPoint[O](180:\rayon){D}
\tkzDefShiftPoint[O](240:\rayon){E}
\tkzDefShiftPoint[O](300:\rayon){F}
\tkzDrawPolygon(A,B,C,D,E,F)
\tkzLabelPoint[right](A){$A$}
\tkzLabelPoint[above right](B){$B$}
\tkzLabelPoint[above left](C){$C$}
\tkzLabelPoint[left](D){$D$}
\tkzLabelPoint[below left](E){$E$}
\tkzLabelPoint[below right](F){$F$}
\tkzDrawSegments[dashed, line width=1pt](O,A O,B O,C O,D O,E O,F) 
\end{tikzpicture} 
}\hfill 
\parbox{0.55\linewidth}{
\begin{enumerate}[label=\alph*)]
\item\rule{0pt}{2.1em} $\vv{EB}+\vv{AF}=$
\item\rule{0pt}{2.1em} $\vv{BC}+\vv{AF}+\vv{DE}=$
\item\rule{0pt}{2.1em} $\vv{OB}+\vv{OE}=$
\item\rule{0pt}{2.1em} $\vv{OA}+\vv{OE}+\vv{OC}=$
\item\rule{0pt}{2.1em} $OA+OE+OC=$
\item\rule{0pt}{2.1em} $\vv{OA}+\vv{OC}+\vv{DE}=$ 
\end{enumerate}
}
\end{exercise}
\vfill 
\begin{exercise}\hfill  

\noindent
\parbox{0.40\linewidth}{\centering
\begin{tikzpicture} 
\tkzSetUpPoint[size=3, fill=black]
\tkzfctset{tan style/.style={->,>=latex, color=red, line width=1.2pt}}   
\tkzSetUpLine[line width= 1.2pt, add=.2 and .2] 

\tkzDefPoint(0,0){O}
\def \rayon{3};  
\tkzDefShiftPoint[O](0:\rayon){A}
\tkzDefShiftPoint[O](60:\rayon){B}
\tkzDefShiftPoint[O](120:\rayon){C}
\tkzDefShiftPoint[O](180:\rayon){D}
\tkzDefShiftPoint[O](240:\rayon){E}
\tkzDefShiftPoint[O](300:\rayon){F}
\tkzDrawPolygon(A,B,C,D,E,F)
\tkzLabelPoint[right](A){$A$}
\tkzLabelPoint[above right](B){$B$}
\tkzLabelPoint[above left](C){$C$}
\tkzLabelPoint[left](D){$D$}
\tkzLabelPoint[below left](E){$E$}
\tkzLabelPoint[below right](F){$F$}
\tkzDrawSegments[dashed, line width=1pt](O,A O,B O,C O,D O,E O,F) 
\end{tikzpicture} 
}\hfill 
\parbox{0.55\linewidth}{
\begin{enumerate}[label=\alph*)] 
\item\rule{0pt}{2.1em} $\vv{OB}-\vv{OA}=$
\item\rule{0pt}{2.1em} $OB-OA=$
\item\rule{0pt}{2.1em} $\vv{CB}-\vv{FA}=$
\item\rule{0pt}{2.1em} $\vv{DO}-\vv{AB}=$
%				\item\rule{0pt}{2.1em} $\vv{DF}-\vv{CB}+\vv{DB}=$
%				\item\rule{0pt}{2.1em} $\vv{OC}-\vv{EC}+\vv{DE}=$
%				\item\rule{0pt}{2.1em} $\vv{OC}-\big(\vv{EO}+\vv{DE}\big)=$
\end{enumerate}
}
\end{exercise}

\begin{exercise}\hfill 

\noindent
\parbox{0.40\linewidth}{\centering
\begin{tikzpicture} 
\tkzSetUpPoint[size=3, fill=black]
\tkzfctset{tan style/.style={->,>=latex, color=red, line width=1.2pt}}   
\tkzSetUpLine[line width= 1.2pt, add=.2 and .2] 

\tkzDefPoint(0,0){O}
\def \rayon{3.0};  
\tkzDefShiftPoint[O](0:\rayon){A}
\tkzDefShiftPoint[O](60:\rayon){B}
\tkzDefShiftPoint[O](120:\rayon){C}
\tkzDefShiftPoint[O](180:\rayon){D}
\tkzDefShiftPoint[O](240:\rayon){E}
\tkzDefShiftPoint[O](300:\rayon){F}
\tkzDrawPolygon(A,B,C,D,E,F)
\tkzLabelPoint[right](A){$A$}
\tkzLabelPoint[above right](B){$B$}
\tkzLabelPoint[above left](C){$C$}
\tkzLabelPoint[left](D){$D$}
\tkzLabelPoint[below left](E){$E$}
\tkzLabelPoint[below right](F){$F$}
\tkzDrawSegments[dashed, line width=1pt](O,A O,B O,C O,D O,E O,F) 
\end{tikzpicture} 
}\hfill 
\parbox{0.55\linewidth}{
\begin{enumerate}[label=\alph*)]
\item\rule{0pt}{2.1em} $\vv{DF}-\vv{CB}+\vv{DB}=$
\item\rule{0pt}{2.1em} $\vv{OC}-\vv{EC}+\vv{DE}=$
\item\rule{0pt}{2.1em} $\vv{OC}-\big(\vv{EO}+\vv{DE}\big)=$
\end{enumerate}
}
\end{exercise}

\begin{example} En utilisant la relation de Chasles, compléter les égalités suivantes où \og $\ldots$\fg{}  représente le nom d'un point :

\begin{enumerate}[label=\alph*)]
\item\rule{0pt}{1.8em} $\vv{AC}=\vv{AB}+\vv{\ldots\phantom{A} C}$ %avant \lstar

\item\rule{0pt}{1.8em} $\vv{FE}=\vv{F\phantom{A}\ldots}+\vv{U\phantom{A}\ldots }$

\item\rule{0pt}{1.8em} $\vv{OU}+\vv{RS}+\vv{UR}=\vv{\ldots\phantom{A}\ldots }$

\item\rule{0pt}{1.8em} $\vv{RT}=\vv{\ldots\phantom{A} I}+\vv{I\phantom{A}\ldots }$

\item\rule{0pt}{1.8em} $\vv{XY}=\vv{\ldots\phantom{A} M}+\vv{\ldots\phantom{A} N}+\vv{\ldots\phantom{A}\ldots }$
\end{enumerate}
\end{example}
\begin{exercise}\label{chap07.exo.chasles01}   
Simplifier les expressions suivantes à l'aide de la relation de Chasles :
\begin{multicols}{2} 
\begin{enumerate} [label=\alph*),leftmargin=*]
\item\rule{0pt}{1.8em} $\vv{BD}+\vv{DA}=\dotfill$
\item\rule{0pt}{1.8em} $\vv{BD}+\vv{DB}=\dotfill$
\item\rule{0pt}{1.8em} $\vv{BD}+\vv{AA}=\dotfill$
\item\rule{0pt}{1.8em} $\vv{BD}-\vv{BA}=\dotfill$
\item\rule{0pt}{1.8em} $\vv{MB}-\vv{MD}=\dotfill$
\item\rule{0pt}{1.8em} $\vv{BD}+\vv{AD}+\vv{BA}=\dotfill$
\item\rule{0pt}{1.8em}  $\vv{KT}+\vv{TD}+\vv{KT}=\dotfill$
\item\rule{0pt}{1.8em}  $\vv{BQ}-\vv{BA}+\vv{QA}-\vv{QB}=\dotfill$
\end{enumerate}
\end{multicols}
\end{exercise}

\begin{exercise}\label{chap07.exo.chasles02}   \hfill 

\noindent
En utilisant la relation de Chasles, compléter les égalités suivantes où \og $\ldots$\fg{}  représente le nom d'un point :
\begin{multicols}{2} 
\begin{enumerate}[label=\alph*),leftmargin=*]
\item\rule{0pt}{1.8em} $\vv{IB}=\vv{\ldots A}+\vv{A\ldots}$
\item\rule{0pt}{1.8em} $\vv{HG}+\vv{\rule[1em]{0cm}{0pt} \ldots\ldots}= \vv{HF}$ 
\item\rule{0pt}{1.8em} $\vv{D\ldots}+\vv{C\ldots}=\vv{\ldots B}$
\item\rule{0pt}{1.8em} $\vv{A\ldots}=\vv{A\ldots}+\vv{B\ldots}+\vv{CM}$
\item\rule{0pt}{1.8em} $\vv{FE}+\vv{\ldots\rule[1em]{0cm}{0pt}\ldots}=\vv{0}$ 

\item\rule{0pt}{1.8em} $\vv{DC} - \vv{\ldots\rule[1em]{0cm}{0pt} \ldots}=\vv{AC}$
\item\rule{0pt}{1.8em} $\vv{DC} - \vv{\ldots\rule[1em]{0cm}{0pt} \ldots}=\vv{DE}$ 
\end{enumerate}
\end{multicols}
\end{exercise}
\begin{exercise}\label{chap07.exo.chasles03} \hfill 

\noindent	Simplifier les expressions suivantes à l'aide de la relation de Chasles :
%\begin{multicols}{2} 
\begin{enumerate} [label=\alph*)] 
\item\rule{0pt}{1.8em} $\vv{BD}-\vv{BA}+\vv{DA}-\vv{DB}=$
\item\rule{0pt}{1.8em} $\vv{CB}-\vv{CD}-\vv{BD}=$
\item\rule{0pt}{1.8em} $\vv{BD}-\vv{BA}+\vv{MA}-\vv{MD}=$
\item\rule{0pt}{1.8em} $\vv{BD}-\vv{MC}-\vv{BM}+\vv{DB}=$
\item\rule{0pt}{1.8em} $\vv{MA}+\vv{EM}-\vv{CA}-\vv{EC}=$
\item\rule{0pt}{1.8em} $-\vv{AU}+\vv{SM}-\vv{ST}+\vv{MU}=$
\end{enumerate}
%\end{multicols}
\end{exercise}

\begin{exercise} \hfill 

\noindent Écrire les sommes demandées à l'aide de vecteurs dont les extrémités sont des points de la figure.


\noindent 
\parbox{0.5\linewidth}{
\begin{enumerate}[label=\alph*)\rule{0pt}{1.9em}, leftmargin=*]
\item $\vv{GF}+\vv{CB}=\dotfill$
\item $\vv{BG}-\vv{HG}=\dotfill$
\item $\vv{HF}+\vv{FH}+\vv{HF}=\dotfill$
\item $\vv{IJ}+\vv{CF}+\vv{JC}+\vv{FE}=\dotfill$
\item $\vv{HF}-\vv{BC}+\vv{CD}=\dotfill$
\item $\vv{BD}+\vv{IH}-\vv{BH}-\vv{FD}=\dotfill$
\end{enumerate}
}\hfill 
\parbox{0.45\linewidth}{\centering
\begin{tikzpicture}[scale=0.75]
\tkzInit[xmax=10,ymin=1, ymax=8]\tkzGrid 
\tkzDefPoint(5,7){A} 
\tkzDefPoint(6,6){B} 
\tkzDefPoint(8,6){C} 
\tkzDefPoint(6,4){D} 
\tkzDefPoint(7,2){E} 
\tkzDefPoint(6,2){F} 
\tkzDefPoint(5,3){G} 
\tkzDefPoint(4,2){H} 
\tkzDefPoint(3,2){I} 
\tkzDefPoint(2,4){J} 
\tkzDrawPolygon[line width=1.2 pt](A,B,C,D,E,F,G,H,I,J)
\tkzLabelPoints(B,C,E,F,H,I)
\tkzLabelPoint[above](A){$A$}
\tkzLabelPoint[left](J){$J$}
\tkzLabelPoint[right](D){$D$}
\tkzLabelPoint[above](G){$G$} 
\end{tikzpicture} 
} 	
\end{exercise}

\begin{multicols}{3}
\begin{proof}[solution de l'exercice \ref{chap07.exo.chasles01}]\hfill 
\begin{enumerate} [label=\alph*), leftmargin=*]\scriptsize
\item $\vv{BA}$
\item $\vv{0}$
\item $\vv{BD}$
\item $\vv{AD}$
\item $\vv{DB}$
\item $2\vv{BD}$
\end{enumerate} 
\end{proof} 
\begin{proof}[solution de  \ref{chap07.exo.chasles02}]\hfill  
\begin{enumerate}[label=\alph*), leftmargin=*]\scriptsize
\item $\vv{IB}=\vv{IA}+\vv{AB}$
\item $\vv{HG}+\vv{GF}=\vv{HF}$
\item $\vv{DC}+\vv{CB}=\vv{DB}$
\item $\vv{AM}=\vv{AB}+\vv{BC}+\vv{CM}$
\item $\vv{FE}+\vv{EF}=\vv{0}$
\item $\vv{DC} - \vv{DA} =\vv{AC}$
\item $\vv{DC} - \vv{EC}  =\vv{DE}$ 
\end{enumerate} 
\end{proof}  
\begin{proof}[solution de  \ref{chap07.exo.chasles03}]\hfill  
\begin{enumerate}[label=\alph*), leftmargin=*]\scriptsize
\item $\vv{BD}$
\item $2\vv{DB}$
\item $\vv{0}$
\item $\vv{CB}$
\item $\vv{0}$
\item $\vv{TA}$
\end{enumerate} 
\end{proof}   
\end{multicols}

\newpage  
\begin{exercise}\hfill\\
Pour chaque cas, représentez $\vv{u}+\vv{v}$ (en bleu) et $\vv{u}-\vv{v}$ (en rouge). Aucune restriction sur l'orgine du représentant choisi.  
%Enfin sur les 3 cas de droite représenter $\vv{v}-\vv{u}$.\\
%	\begin{minipage}{1\columnwidth}\centering
%		\begin{tikzpicture}[scale=0.5]
%			\tkzInit[xmax=34,ymax=30]
%			%	\begin{scope}[dashed]
%			\tkzGrid 
%			%	\end{scope}
%			\tkzSetUpPoint[size=7, fill=black]
%			\tkzfctset{tan style/.style={->,>=latex, color=red, line width=1.2pt}} 
%			\tkzDefPoint(1,25){P1}
%			\tkzDefShiftPoint[P1](6,3){Q1}
%			\tkzDefShiftPoint[Q1](3,-2){R1}
%			\tkzDrawSegment[->,>=latex, line width=1.2](P1,Q1) 
%			\tkzDrawSegment[->,>=latex, line width=1.2](Q1,R1) 
%			\tkzLabelSegment[above left](P1,Q1){$\vv{u}$}
%			\tkzLabelSegment[above right](Q1,R1){$\vv{v}$}	
%			
%			\tkzDefPoint(15,26){P2}
%			\tkzDefShiftPoint[P2](3,2){Q2}
%			\tkzDefShiftPoint[Q2](5,-1){R2}
%			\tkzDrawSegment[->,>=latex, line width=1.2](Q2,P2) 
%			\tkzDrawSegment[->,>=latex, line width=1.2](Q2,R2) 
%			\tkzLabelSegment[above left](P2,Q2){$\vv{u}$}
%			\tkzLabelSegment[above right](Q2,R2){$\vv{v}$}		
%			
%			\tkzDefPoint(33,29){P3}
%			\tkzDefShiftPoint[P3](-6,0){Q3}
%			\tkzDefShiftPoint[Q3](0,-4){R3}
%			\tkzDrawSegment[->,>=latex, line width=1.2](P3,Q3) 
%			\tkzDrawSegment[->,>=latex, line width=1.2](Q3,R3) 
%			\tkzLabelSegment[above left](P3,Q3){$\vv{v}$}
%			\tkzLabelSegment[above right](Q3,R3){$\vv{u}$}	
%			
%			\tkzDefPoint(1,15){P4}
%			\tkzDefShiftPoint[P4](6,3){Q4}
%			\tkzDefPoint(7,16){R4}
%			\tkzDefShiftPoint[R4](4,-1){S4}
%			\tkzDrawSegment[->,>=latex, line width=1.2](P4,Q4) 
%			\tkzDrawSegment[->,>=latex, line width=1.2](R4,S4) 
%			\tkzLabelSegment[above left](P4,Q4){$\vv{u}$}
%			\tkzLabelSegment[above right](R4,S4){$\vv{v}$}	
%			
%			\tkzDefPoint(18,17){P5}
%			\tkzDefShiftPoint[P5](-4,0){Q5}
%			\tkzDefPoint(22,12){R5}
%			\tkzDefShiftPoint[R5](-2,5){S5}
%			\tkzDrawSegment[->,>=latex, line width=1.2](P5,Q5) 
%			\tkzDrawSegment[->,>=latex, line width=1.2](R5,S5) 
%			\tkzLabelSegment[above left](P5,Q5){$\vv{u}$}
%			\tkzLabelSegment[above right](R5,S5){$\vv{v}$}	
%			
%			\tkzDefPoint(1, 3){P6}
%			\tkzDefShiftPoint[P6](5,0){Q6}
%			\tkzDefPoint(8,6){R6}
%			\tkzDefShiftPoint[R6](-2,-3){S6}
%			\tkzDrawSegment[->,>=latex, line width=1.2](P6,Q6) 
%			\tkzDrawSegment[->,>=latex, line width=1.2](R6,S6) 
%			\tkzLabelSegment[above left](P6,Q6){$\vv{u}$}
%			\tkzLabelSegment[below right](R6,S6){$\vv{v}$}	
%			
%			\tkzDefPoint(33, 16){P7}
%			\tkzDefShiftPoint[P7](-6,0){Q7}
%			\tkzDefPoint(33, 13){R7}
%			\tkzDefShiftPoint[R7](-4,0){S7}
%			\tkzDrawSegment[->,>=latex, line width=1.2](P7,Q7) 
%			\tkzDrawSegment[->,>=latex, line width=1.2](R7,S7) 
%			\tkzLabelSegment[above left](P7,Q7){$\vv{u}$}
%			\tkzLabelSegment[below right](R7,S7){$\vv{v}$}
%			
%			
%			\tkzDefPoint(20, 2){P8}
%			\tkzDefShiftPoint[P8](-4,0){Q8}
%			\tkzDefPoint(14, 4){R8}
%			\tkzDefShiftPoint[R8](6,0){S8}
%			\tkzDrawSegment[->,>=latex, line width=1.2](P8,Q8) 
%			\tkzDrawSegment[->,>=latex, line width=1.2](R8,S8) 
%			\tkzLabelSegment[ below](P8,Q8){$\vv{u}$}
%			\tkzLabelSegment[above  ](R8,S8){$\vv{v}$}
%			
%			
%			\tkzDefPoint(31, 3){P9}
%			\tkzDefShiftPoint[P9](-3,0){Q9}
%			\tkzDefPoint(25, 3){R9}
%			\tkzDefShiftPoint[R9](3,0){S9}
%			\tkzDrawSegment[->,>=latex, line width=1.2](P9,Q9) 
%			\tkzDrawSegment[->,>=latex, line width=1.2](R9,S9) 
%			\tkzLabelSegment[ below](P9,Q9){$\vv{u}$}
%			\tkzLabelSegment[above  ](R9,S9){$\vv{v}$}
%			
%			%		\tkzAxeX[orig=false,above=5pt, label=$x$, font=\scriptsize, line width=2pt, right space=0.2,5 cm] 
%			%		\tkzAxeY[orig=false,right=5pt, label=$y$, font=\scriptsize, line width=2pt, up space=0.2,5 cm]
%		\end{tikzpicture} 
%	\end{minipage}  
\begin{figure}[H]
\begin{tikzpicture}[scale=0.5]
\tkzInit[xmax=34,ymax=42]
%	\begin{scope}[dashed]
\tkzGrid 
%	\end{scope}
\tkzSetUpPoint[size=7, fill=black]
\tkzfctset{tan style/.style={->,>=latex, color=red, line width=1.2pt}} 
\tkzDefPoint(1,37){P1}
\tkzDefShiftPoint[P1](6,3){Q1}
\tkzDefShiftPoint[Q1](3,-2){R1}
\tkzDrawSegment[->,>=latex, line width=1.2](P1,Q1) 
\tkzDrawSegment[->,>=latex, line width=1.2](Q1,R1) 
\tkzLabelSegment[above left](P1,Q1){$\vv{u}$}
\tkzLabelSegment[above right](Q1,R1){$\vv{v}$}	

\tkzDefPoint(15,38){P2}
\tkzDefShiftPoint[P2](3,2){Q2}
\tkzDefShiftPoint[Q2](5,-1){R2}
\tkzDrawSegment[->,>=latex, line width=1.2](Q2,P2) 
\tkzDrawSegment[->,>=latex, line width=1.2](Q2,R2) 
\tkzLabelSegment[above left](P2,Q2){$\vv{u}$}
\tkzLabelSegment[above right](Q2,R2){$\vv{v}$}		

\tkzDefPoint(33,41){P3}
\tkzDefShiftPoint[P3](-6,0){Q3}
\tkzDefShiftPoint[Q3](0,-4){R3}
\tkzDrawSegment[->,>=latex, line width=1.2](P3,Q3) 
\tkzDrawSegment[->,>=latex, line width=1.2](Q3,R3) 
\tkzLabelSegment[above left](P3,Q3){$\vv{v}$}
\tkzLabelSegment[above right](Q3,R3){$\vv{u}$}	

\tkzDefPoint(1,21){P4}
\tkzDefShiftPoint[P4](6,3){Q4}
\tkzDefPoint(7,22){R4}
\tkzDefShiftPoint[R4](4,-1){S4}
\tkzDrawSegment[->,>=latex, line width=1.2](P4,Q4) 
\tkzDrawSegment[->,>=latex, line width=1.2](R4,S4) 
\tkzLabelSegment[above left](P4,Q4){$\vv{u}$}
\tkzLabelSegment[above right](R4,S4){$\vv{v}$}	

\tkzDefPoint(18,23){P5}
\tkzDefShiftPoint[P5](-4,0){Q5}
\tkzDefPoint(22,18){R5}
\tkzDefShiftPoint[R5](-2,5){S5}
\tkzDrawSegment[->,>=latex, line width=1.2](P5,Q5) 
\tkzDrawSegment[->,>=latex, line width=1.2](R5,S5) 
\tkzLabelSegment[above left](P5,Q5){$\vv{u}$}
\tkzLabelSegment[above right](R5,S5){$\vv{v}$}	

\tkzDefPoint(1, 3){P6}
\tkzDefShiftPoint[P6](5,0){Q6}
\tkzDefPoint(8,6){R6}
\tkzDefShiftPoint[R6](-2,-3){S6}
\tkzDrawSegment[->,>=latex, line width=1.2](P6,Q6) 
\tkzDrawSegment[->,>=latex, line width=1.2](R6,S6) 
\tkzLabelSegment[above left](P6,Q6){$\vv{u}$}
\tkzLabelSegment[below right](R6,S6){$\vv{v}$}	

\tkzDefPoint(33, 22){P7}
\tkzDefShiftPoint[P7](-6,0){Q7}
\tkzDefPoint(33, 19){R7}
\tkzDefShiftPoint[R7](-4,0){S7}
\tkzDrawSegment[->,>=latex, line width=1.2](P7,Q7) 
\tkzDrawSegment[->,>=latex, line width=1.2](R7,S7) 
\tkzLabelSegment[above left](P7,Q7){$\vv{u}$}
\tkzLabelSegment[below right](R7,S7){$\vv{v}$}


\tkzDefPoint(20, 2){P8}
\tkzDefShiftPoint[P8](-4,0){Q8}
\tkzDefPoint(14, 4){R8}
\tkzDefShiftPoint[R8](6,0){S8}
\tkzDrawSegment[->,>=latex, line width=1.2](P8,Q8) 
\tkzDrawSegment[->,>=latex, line width=1.2](R8,S8) 
\tkzLabelSegment[ below](P8,Q8){$\vv{u}$}
\tkzLabelSegment[above  ](R8,S8){$\vv{v}$}


\tkzDefPoint(31, 3){P9}
\tkzDefShiftPoint[P9](-3,0){Q9}
\tkzDefPoint(25, 3){R9}
\tkzDefShiftPoint[R9](3,0){S9}
\tkzDrawSegment[->,>=latex, line width=1.2](P9,Q9) 
\tkzDrawSegment[->,>=latex, line width=1.2](R9,S9) 
\tkzLabelSegment[ below](P9,Q9){$\vv{u}$}
\tkzLabelSegment[above  ](R9,S9){$\vv{v}$}

%		\tkzAxeX[orig=false,above=5pt, label=$x$, font=\scriptsize, line width=2pt, right space=0.2,5 cm] 
%		\tkzAxeY[orig=false,right=5pt, label=$y$, font=\scriptsize, line width=2pt, up space=0.2,5 cm]
\end{tikzpicture} 
\end{figure} 
\end{exercise}  

%%%%%%%%%%%%%%%%%%%%%%%%%%%%%%%%%%%%%%%%%%%%%%%%%%%%%%%%%%%
%
%%%%%%%%%%%%%%%%%%%%%%%%%%%%%%%%%%%%%%%%%%%%%%%%%%%%%%%%%%%
\newpage 
\begin{exercise}\hfill \\
Pour chacun des cas suivants placer le point $M$ tel que $\vv{OM}=\vv{OA}+\vv{OB}$ et représenter le parallélogramme $OAMB$ puis écrire le vecteur $\vv{AB}$ sous la forme $a\vv{\imath}+b\vv{\jmath}$. 
\begin{multicols}{3} 
\begin{enumerate}[label=\alph*)]
\item $A (1;4)$ et $B  (3;2)$.
\item $A (-1;2)$ et $B  (3;1)$. 
\item $A (-1;1)$ et $B  (-1;2)$. 
\item $A (1;5)$ et $B  (1;1)$. 
\item $A (-2;1)$ et $B  (1;2)$. 
\item $A (-3;-2)$ et $B  (-1;-1)$.  
\end{enumerate} 
\end{multicols}
\end{exercise}
%\end{minipage}

\begin{minipage}{0.45\columnwidth} \centering
\begin{tikzpicture}[xscale=0.75, yscale=0.65]

\tkzInit[xmin=-4, ymin=-3,  xmax=6,ymax=6 ]
\begin{scope}[dashed] 
\tkzGrid[ xstep=1, ystep=1](-4,-3)(6,6)
\end{scope}

\def \xa{1}; 
\def \ya{4}; 

\def \xb{3};
\def \yb{2};

\tkzDefPoint(0,0){O} \tkzDefShiftPoint[O](\xa,\ya){A}
\tkzDefShiftPoint[O](\xb,\yb){B}
\tkzDrawSegments[- ,>=latex,line width =1.2 pt, color = Maroon](O,A O,B)
%\tkzLabelSegment[above](U0,U1){$\vv{u}$}
%\tkzLabelSegment[left](V0,V1){$\vv{v}$}

\tkzDrawX[left=5pt, label=, font=\scriptsize, line width=1.2pt ] %, style=dashed
\tkzDrawY[below =5pt, label=, font=\scriptsize, line width=1.2pt]  %, style=dashed
\tkzLabelXY[orig = false,  font=\scriptsize]

\tkzDrawPoints(A,B,O)
\tkzLabelPoints[below left, font=\scriptsize](O)
\tkzLabelPoints[font=\scriptsize](A,B)  
\tkzRep[line width=2pt, color=Maroon]
\end{tikzpicture} 
\end{minipage} 
\hfill 
\begin{minipage}{0.45\columnwidth} \centering
\begin{tikzpicture}[xscale=0.75, yscale=0.65]

\tkzInit[xmin=-4, ymin=-3,  xmax=6,ymax=6 ]
\begin{scope}[dashed] 
\tkzGrid[ xstep=1, ystep=1](-4,-3)(6,6)
\end{scope}

\def \xa{1}; 
\def \ya{4}; 

\def \xb{3};
\def \yb{2};

%\tkzDefPoint(0,0){O} \tkzDefShiftPoint[O](\xa,\ya){A}
% \tkzDefShiftPoint[O](\xb,\yb){B}
%\tkzDrawSegments[- ,>=latex,line width =2 pt, color = Maroon](O,A O,B)
%\tkzLabelSegment[above](U0,U1){$\vv{u}$}
%\tkzLabelSegment[left](V0,V1){$\vv{v}$}

\tkzDrawX[left=5pt, label=, font=\scriptsize, line width=1.2pt ] %, style=dashed
\tkzDrawY[below =5pt, label=, font=\scriptsize, line width=1.2pt]  %, style=dashed
\tkzLabelXY[orig = false,  font=\scriptsize]

\tkzLabelPoints[below left, font=\scriptsize](O) 
\tkzRep[line width=2pt, color=Maroon]
\end{tikzpicture} 
\end{minipage} 

\begin{minipage}{0.45\columnwidth} \centering
\begin{tikzpicture}[xscale=0.75, yscale=0.65]

\tkzInit[xmin=-4, ymin=-3,  xmax=6,ymax=6 ]
\begin{scope}[dashed] 
\tkzGrid[ xstep=1, ystep=1](-4,-3)(6,6)
\end{scope}

\def \xa{1}; 
\def \ya{4}; 

\def \xb{3};
\def \yb{2};

%\tkzDefPoint(0,0){O} \tkzDefShiftPoint[O](\xa,\ya){A}
% \tkzDefShiftPoint[O](\xb,\yb){B}
%\tkzDrawSegments[- ,>=latex,line width =2 pt, color = Maroon](O,A O,B)
%\tkzLabelSegment[above](U0,U1){$\vv{u}$}
%\tkzLabelSegment[left](V0,V1){$\vv{v}$}

\tkzDrawX[left=5pt, label=, font=\scriptsize, line width=1.2pt ] %, style=dashed
\tkzDrawY[below =5pt, label=, font=\scriptsize, line width=1.2pt]  %, style=dashed
\tkzLabelXY[orig = false,  font=\scriptsize]

\tkzLabelPoints[below left, font=\scriptsize](O)
\tkzRep[line width=2pt, color=Maroon]
\end{tikzpicture} 
\end{minipage} 
\hfill 
\begin{minipage}{0.45\columnwidth} \centering
\begin{tikzpicture}[xscale=0.75, yscale=0.65]

\tkzInit[xmin=-4, ymin=-3,  xmax=6,ymax=6 ]
\begin{scope}[dashed] 
\tkzGrid[ xstep=1, ystep=1](-4,-3)(6,6)
\end{scope}

\def \xa{1}; 
\def \ya{4}; 

\def \xb{3};
\def \yb{2};

%\tkzDefPoint(0,0){O} \tkzDefShiftPoint[O](\xa,\ya){A}
% \tkzDefShiftPoint[O](\xb,\yb){B}
%\tkzDrawSegments[- ,>=latex,line width =2 pt, color = Maroon](O,A O,B)
%\tkzLabelSegment[above](U0,U1){$\vv{u}$}
%\tkzLabelSegment[left](V0,V1){$\vv{v}$}

\tkzDrawX[left=5pt, label=, font=\scriptsize, line width=1.2pt ] %, style=dashed
\tkzDrawY[below =5pt, label=, font=\scriptsize, line width=1.2pt]  %, style=dashed
\tkzLabelXY[orig = false,  font=\scriptsize]

\tkzLabelPoints[below left, font=\scriptsize](O) 
\tkzRep[line width=2pt, color=Maroon]
\end{tikzpicture} 
\end{minipage} 


\begin{minipage}{0.45\columnwidth} \centering
\begin{tikzpicture}[xscale=0.75, yscale=0.65]

\tkzInit[xmin=-4, ymin=-3,  xmax=6,ymax=6 ]
\begin{scope}[dashed] 
\tkzGrid[ xstep=1, ystep=1](-4,-3)(6,6)
\end{scope}

\def \xa{1}; 
\def \ya{4}; 

\def \xb{3};
\def \yb{2};

%\tkzDefPoint(0,0){O} \tkzDefShiftPoint[O](\xa,\ya){A}
% \tkzDefShiftPoint[O](\xb,\yb){B}
%\tkzDrawSegments[- ,>=latex,line width =2 pt, color = Maroon](O,A O,B)
%\tkzLabelSegment[above](U0,U1){$\vv{u}$}
%\tkzLabelSegment[left](V0,V1){$\vv{v}$}

\tkzDrawX[left=5pt, label=, font=\scriptsize, line width=1.2pt ] %, style=dashed
\tkzDrawY[below =5pt, label=, font=\scriptsize, line width=1.2pt]  %, style=dashed
\tkzLabelXY[orig = false,  font=\scriptsize]

\tkzLabelPoints[below left, font=\scriptsize](O)
\tkzRep[line width=2pt, color=Maroon]
\end{tikzpicture} 
\end{minipage} 
\hfill 
\begin{minipage}{0.45\columnwidth} \centering
\begin{tikzpicture}[xscale=0.75, yscale=0.65]

\tkzInit[xmin=-4, ymin=-3,  xmax=6,ymax=6 ]
\begin{scope}[dashed] 
\tkzGrid[ xstep=1, ystep=1](-4,-3)(6,6)
\end{scope}

\def \xa{1}; 
\def \ya{4}; 

\def \xb{3};
\def \yb{2};

%\tkzDefPoint(0,0){O} \tkzDefShiftPoint[O](\xa,\ya){A}
% \tkzDefShiftPoint[O](\xb,\yb){B}
%\tkzDrawSegments[- ,>=latex,line width =2 pt, color = Maroon](O,A O,B)
%\tkzLabelSegment[above](U0,U1){$\vv{u}$}
%\tkzLabelSegment[left](V0,V1){$\vv{v}$}

\tkzDrawX[left=5pt, label=, font=\scriptsize, line width=1.2pt ] %, style=dashed
\tkzDrawY[below =5pt, label=, font=\scriptsize, line width=1.2pt]  %, style=dashed
\tkzLabelXY[orig = false,  font=\scriptsize]

\tkzLabelPoints[below left, font=\scriptsize](O) 
\tkzRep[line width=2pt, color=Maroon]
\end{tikzpicture} 
\end{minipage} 


%%%%%%%%%%%%%%%%%%%%%%%%%%%%%%%%%%%%%%%%%%%%%%%%%%%%%%%%%%%
%
%%%%%%%%%%%%%%%%%%%%%%%%%%%%%%%%%%%%%%%%%%%%%%%%%%%%%%%%%%%
\newpage 
\begin{exercise}[décomposer un vecteur en fonction de deux autres]\label{chapter16:vecteur:decompositionA}\hfill 

\noindent La figure ci-dessous (répétée 3 fois) est formée de triangles équilatéraux. On pose $\vv{OC}=\vv{\imath}$ et $\vv{OD}=\vv{\jmath}$. Écrire chacun des vecteurs suivants sous la forme $a\vv{\imath}+b\vv{\jmath}$.

\noindent
\parbox{0.32\linewidth}{
\begin{tikzpicture} 
\tkzSetUpPoint[size=3, fill=black]
\tkzfctset{tan style/.style={->,>=latex, color=red, line width=1.2pt}}   
\tkzSetUpLine[line width= 1.2pt, add=.2 and .2] 

\tkzDefPoint(0,0){O}
\def \rayon{1.8};  
\tkzDefShiftPoint[O](0:\rayon){C}
\tkzDefShiftPoint[O](60:\rayon){D}
\tkzDefShiftPoint[O](120:\rayon){E}
\tkzDefShiftPoint[O](180:\rayon){F}
\tkzDefShiftPoint[O](240:\rayon){A}
\tkzDefShiftPoint[O](300:\rayon){B}
\tkzDefShiftPoint[B](0:\rayon){H}
\tkzDefShiftPoint[E](60:\rayon){I}
\tkzDefShiftPoint[A](180:\rayon){G}
%	\tkzDrawPolygon(A,B,C,D,E,F,G,H,I)
\tkzLabelPoint[above](O){\scriptsize$O$}
\tkzLabelPoint[below](A){\scriptsize$A$}
\tkzLabelPoint[below](B){\scriptsize$B$}
\tkzLabelPoint[above right](C){\scriptsize$C$}
\tkzLabelPoint[above right](D){\scriptsize$D$}
\tkzLabelPoint[above left](E){\scriptsize$E$}
\tkzLabelPoint[above left](F){\scriptsize$F$} 
\tkzLabelPoint[below](G){\scriptsize$G$}
\tkzLabelPoint[below](H){\scriptsize$H$}
\tkzLabelPoint[above ](I){\scriptsize$I$}
\tkzDrawSegments[dashed, line width=1pt](G,H G,I I,H E,B F,A A,D B,C F,C E,D) 
\tkzDrawSegments[->, >=latex,color=red,line width=1.2pt](O,C O,D)
\tkzLabelSegment[below](O,C){$\vv{\imath}$}
\tkzLabelSegment[below](O,D){$\vv{\jmath}$} 
\end{tikzpicture}
}
\hfill
\parbox{0.32\linewidth}{
\begin{tikzpicture} 
\tkzSetUpPoint[size=3, fill=black]
\tkzfctset{tan style/.style={->,>=latex, color=red, line width=1.2pt}}   
\tkzSetUpLine[line width= 1.2pt, add=.2 and .2] 

\tkzDefPoint(0,0){O}
\def \rayon{1.8};  
\tkzDefShiftPoint[O](0:\rayon){C}
\tkzDefShiftPoint[O](60:\rayon){D}
\tkzDefShiftPoint[O](120:\rayon){E}
\tkzDefShiftPoint[O](180:\rayon){F}
\tkzDefShiftPoint[O](240:\rayon){A}
\tkzDefShiftPoint[O](300:\rayon){B}
\tkzDefShiftPoint[B](0:\rayon){H}
\tkzDefShiftPoint[E](60:\rayon){I}
\tkzDefShiftPoint[A](180:\rayon){G}
%	\tkzDrawPolygon(A,B,C,D,E,F,G,H,I)
\tkzLabelPoint[above](O){\scriptsize$O$}
\tkzLabelPoint[below](A){\scriptsize$A$}
\tkzLabelPoint[below](B){\scriptsize$B$}
\tkzLabelPoint[above right](C){\scriptsize$C$}
\tkzLabelPoint[above right](D){\scriptsize$D$}
\tkzLabelPoint[above left](E){\scriptsize$E$}
\tkzLabelPoint[above left](F){\scriptsize$F$} 
\tkzLabelPoint[below](G){\scriptsize$G$}
\tkzLabelPoint[below](H){\scriptsize$H$}
\tkzLabelPoint[above ](I){\scriptsize$I$}
\tkzDrawSegments[dashed, line width=1pt](G,H G,I I,H E,B F,A A,D B,C F,C E,D) 
\tkzDrawSegments[->, >=latex,color=red,line width=1.2pt](O,C O,D)
\tkzLabelSegment[below](O,C){$\vv{\imath}$}
\tkzLabelSegment[below](O,D){$\vv{\jmath}$} 
\end{tikzpicture}
}\hfill
\parbox{0.32\linewidth}{
\begin{tikzpicture} 
\tkzSetUpPoint[size=3, fill=black]
\tkzfctset{tan style/.style={->,>=latex, color=red, line width=1.2pt}}   
\tkzSetUpLine[line width= 1.2pt, add=.2 and .2] 

\tkzDefPoint(0,0){O}
\def \rayon{1.8};  
\tkzDefShiftPoint[O](0:\rayon){C}
\tkzDefShiftPoint[O](60:\rayon){D}
\tkzDefShiftPoint[O](120:\rayon){E}
\tkzDefShiftPoint[O](180:\rayon){F}
\tkzDefShiftPoint[O](240:\rayon){A}
\tkzDefShiftPoint[O](300:\rayon){B}
\tkzDefShiftPoint[B](0:\rayon){H}
\tkzDefShiftPoint[E](60:\rayon){I}
\tkzDefShiftPoint[A](180:\rayon){G}
%	\tkzDrawPolygon(A,B,C,D,E,F,G,H,I)
\tkzLabelPoint[above](O){\scriptsize$O$}
\tkzLabelPoint[below](A){\scriptsize$A$}
\tkzLabelPoint[below](B){\scriptsize$B$}
\tkzLabelPoint[above right](C){\scriptsize$C$}
\tkzLabelPoint[above right](D){\scriptsize$D$}
\tkzLabelPoint[above left](E){\scriptsize$E$}
\tkzLabelPoint[above left](F){\scriptsize$F$} 
\tkzLabelPoint[below](G){\scriptsize$G$}
\tkzLabelPoint[below](H){\scriptsize$H$}
\tkzLabelPoint[above ](I){\scriptsize$I$}
\tkzDrawSegments[dashed, line width=1pt](G,H G,I I,H E,B F,A A,D B,C F,C E,D) 
\tkzDrawSegments[->, >=latex,color=red,line width=1.2pt](O,C O,D)
\tkzLabelSegment[below](O,C){$\vv{\imath}$}
\tkzLabelSegment[below](O,D){$\vv{\jmath}$} 
\end{tikzpicture}
}
\begin{multicols}{3}
\begin{enumerate}[label={\alph*)\rule{0pt}{2.1em}},leftmargin=*] 
\item $\vv{BI}=$ \phantom{$-2\vv{\imath}+3\vv{\jmath}$} 
\item $\vv{IA}=$ \phantom{$\vv{\imath}-3\vv{\jmath}$} 
\item $\vv{IH}=$ \phantom{$3\vv{\imath}-3\vv{\jmath}$} 
\item $\vv{EB}=$ \phantom{$2\vv{\imath}-2\vv{\jmath}$} 
\item $\vv{OG}=$ \phantom{$-\vv{\imath}-\vv{\jmath}$} 
\item $\vv{DG}=$ \phantom{$-\vv{\imath}-2\vv{\jmath}$} 
\item $\vv{GC}=$ \phantom{$2\vv{\imath}+\vv{\jmath}$} 
\item $\vv{HE}=$ \phantom{$-3\vv{\imath}+2\vv{\jmath}$} 
\item $\vv{AE}=$ \phantom{$-\vv{\imath}+2\vv{\jmath}$} 
\item $\vv{AF}=$ \phantom{$-\vv{\imath}+\vv{\jmath}$} 
\item $\vv{BD}=$ \phantom{$-\vv{\imath}+2\vv{\jmath}$} 
\item $\vv{FB}=$ \phantom{$2\vv{\imath}-\vv{\jmath}$}  
\end{enumerate}
\end{multicols}




\end{exercise}
\begin{exercise}[décomposer un vecteur en fonction de deux autres]\label{chapter16:vecteur:decompositionB}  La figure ci-dessous est formée de deux hexagones réguliers de centre $O$. $I$, $J$, $K$, $L$, $M$ et $N$ sont les milieux respectifs des segements $[OA]$, $[OB]$, $[OC]$, $[OD]$, $[OE]$ et $[OF]$. 


\noindent
\parbox{0.32\linewidth}{	
\begin{tikzpicture} 
\tkzSetUpPoint[size=3, fill=black]
\tkzfctset{tan style/.style={->,>=latex, color=red, line width=1.2pt}}   
\tkzSetUpLine[line width= 1.2pt, add=.2 and .2] 

\tkzDefPoint(0,0){O}
\def \rayon{2.3};  
\tkzDefShiftPoint[O](-120:\rayon){A}
\tkzDefShiftPoint[O](-60:\rayon){B}
\tkzDefShiftPoint[O](0:\rayon){C}
\tkzDefShiftPoint[O](60:\rayon){D}
\tkzDefShiftPoint[O](120:\rayon){E}
\tkzDefShiftPoint[O](180:\rayon){F}
\tkzDrawPoints(O,A,B,C,D,E,F)
\tkzDrawPolygon(A,B,C,D,E,F)  
\tkzLabelPoint[below](O){\scriptsize$O$}
\tkzLabelPoint[below](A){\scriptsize$A$}
\tkzLabelPoint[below](B){\scriptsize$B$}
\tkzLabelPoint[right](C){\scriptsize$C$}
\tkzLabelPoint[above](D){\scriptsize$D$}
\tkzLabelPoint[above](E){\scriptsize$E$}
\tkzLabelPoint[left](F){\scriptsize$F$}
\tkzDefShiftPoint[O](0:{0.5*\rayon}){K}
\tkzDefShiftPoint[O](60:{0.5*\rayon}){L}
\tkzDefShiftPoint[O](120:{0.5*\rayon}){M}
\tkzDefShiftPoint[O](180:{0.5*\rayon}){N}
\tkzDefShiftPoint[O](240:{0.5*\rayon}){I}
\tkzDefShiftPoint[O](300:{0.5*\rayon}){J}
\tkzDrawPolygon(I,J,K,L,M,N)
\tkzLabelPoint[below](I){\scriptsize$I$}
\tkzLabelPoint[below](J){\scriptsize$J$}
\tkzLabelPoint[above](M){\scriptsize$M$}
\tkzLabelPoint[above](L){\scriptsize$L$}
\tkzLabelPoint[above right](K){\scriptsize$K$}
\tkzLabelPoint[above left](N){\scriptsize$N$}
\tkzDrawSegments[dashed, line width=1pt](O,A O,B O,C O,D O,E O,F)  
%		\tkzDrawSegments[->, >=latex, line width=1.2pt](O,I O,J)
%		\tkzLabelSegment[below](O,I){$\vv{\imath}$}
%		\tkzLabelSegment[below](O,J){$\vv{\jmath}$} 
\end{tikzpicture}	
} \hfill 
\parbox{0.65\linewidth}{
Décomposer chacun des vecteurs suivants selon les vecteurs donnés.
\begin{multicols}{2}
\begin{enumerate}[label={\alph*)\rule{0pt}{2.1em}},leftmargin=*] 
\item $\vv{BF}=\ldots \vv{OI}+\ldots\vv{OJ}$ %\phantom{$2\vv{OI}-4\vv{OJ}$} 
\item $\vv{JD}=\ldots \vv{OI}+\ldots\vv{OK}$ %\phantom{$-3\vv{OI}-\vv{OK}$} 
\item $\vv{CO}=\ldots \vv{OI}+\ldots\vv{OJ}$ %\phantom{$2\vv{OI}-2\vv{OJ}$} 
\item $\vv{EA}=\ldots \vv{OI}+\ldots\vv{OK}$ %\phantom{$4\vv{OI}+2\vv{OK}$}  
\end{enumerate}
\end{multicols}
}


\noindent
\parbox{0.32\linewidth}{	
\begin{tikzpicture} 
\tkzSetUpPoint[size=3, fill=black]
\tkzfctset{tan style/.style={->,>=latex, color=red, line width=1.2pt}}   
\tkzSetUpLine[line width= 1.2pt, add=.2 and .2] 

\tkzDefPoint(0,0){O}
\def \rayon{2.3}; 
\tkzDefShiftPoint[O](-120:\rayon){A}
\tkzDefShiftPoint[O](-60:\rayon){B}
\tkzDefShiftPoint[O](0:\rayon){C}
\tkzDefShiftPoint[O](60:\rayon){D}
\tkzDefShiftPoint[O](120:\rayon){E}
\tkzDefShiftPoint[O](180:\rayon){F}
\tkzDrawPoints(O,A,B,C,D,E,F)
\tkzDrawPolygon(A,B,C,D,E,F)  
\tkzLabelPoint[below](O){\scriptsize$O$}
\tkzLabelPoint[below](A){\scriptsize$A$}
\tkzLabelPoint[below](B){\scriptsize$B$}
\tkzLabelPoint[right](C){\scriptsize$C$}
\tkzLabelPoint[above](D){\scriptsize$D$}
\tkzLabelPoint[above](E){\scriptsize$E$}
\tkzLabelPoint[left](F){\scriptsize$F$}
\tkzDefShiftPoint[O](0:{0.5*\rayon}){K}
\tkzDefShiftPoint[O](60:{0.5*\rayon}){L}
\tkzDefShiftPoint[O](120:{0.5*\rayon}){M}
\tkzDefShiftPoint[O](180:{0.5*\rayon}){N}
\tkzDefShiftPoint[O](240:{0.5*\rayon}){I}
\tkzDefShiftPoint[O](300:{0.5*\rayon}){J}
\tkzDrawPolygon(I,J,K,L,M,N)
\tkzLabelPoint[below](I){\scriptsize$I$}
\tkzLabelPoint[below](J){\scriptsize$J$}
\tkzLabelPoint[above](M){\scriptsize$M$}
\tkzLabelPoint[above](L){\scriptsize$L$}
\tkzLabelPoint[above right](K){\scriptsize$K$}
\tkzLabelPoint[above left](N){\scriptsize$N$}
\tkzDrawSegments[dashed, line width=1pt](O,A O,B O,C O,D O,E O,F)  
%		\tkzDrawSegments[->, >=latex, line width=1.2pt](O,A O,B)
%		\tkzLabelSegment[below](O,I){$\vv{\imath}$}
%		\tkzLabelSegment[below](O,J){$\vv{\jmath}$} 
\end{tikzpicture}	
} \hfill 
\parbox{0.65\linewidth}{
Écrire chacun des vecteurs suivants sous la forme $a\vv{OA}+b\vv{OB}$.
\begin{multicols}{2}
\begin{enumerate}[label={\alph*)\rule{0pt}{2.1em}},leftmargin=*] 
\item $\vv{JF}=\ldots \vv{OA}+\ldots\vv{OB}$ %\phantom{$\vv{OA}-\tfrac{3}{2}\vv{OB}$} 
\item $\vv{EK}=\ldots \vv{OA}+\ldots\vv{OC}$ %\phantom{$\vv{OA}+\tfrac{3}{2}\vv{OC}$} 
\item $\vv{CI}=\ldots \vv{OA}+\ldots\vv{OB}$ %\phantom{$\tfrac{3}{2}\vv{OA}-\vv{OB}$} 
\item $\vv{JD}=\ldots \vv{OA}+\ldots\vv{OC}$ %\phantom{$-\tfrac{3}{2}\vv{OA}-\tfrac{1}{2}\vv{OC}$}   
\end{enumerate}
\end{multicols}
} 
\end{exercise}
%----------------------------------------------------------------------------------------
%	 
%----------------------------------------------------------------------------------------

\newpage 
\begin{proof}[solution de \ref{chapter16:vecteur:decompositionA}]\hfill 

\noindent
\parbox{0.32\linewidth}{
\begin{tikzpicture} 
\tkzSetUpPoint[size=3, fill=black]
\tkzfctset{tan style/.style={->,>=latex, color=red, line width=1.2pt}}   
\tkzSetUpLine[line width= 1.2pt, add=.2 and .2] 

\tkzDefPoint(0,0){O}
\def \rayon{1.8};  
\tkzDefShiftPoint[O](0:\rayon){C}
\tkzDefShiftPoint[O](60:\rayon){D}
\tkzDefShiftPoint[O](120:\rayon){E}
\tkzDefShiftPoint[O](180:\rayon){F}
\tkzDefShiftPoint[O](240:\rayon){A}
\tkzDefShiftPoint[O](300:\rayon){B}
\tkzDefShiftPoint[B](0:\rayon){H}
\tkzDefShiftPoint[E](60:\rayon){I}
\tkzDefShiftPoint[A](180:\rayon){G}
%	\tkzDrawPolygon(A,B,C,D,E,F,G,H,I)
\tkzLabelPoint[above](O){\scriptsize$O$}
\tkzLabelPoint[below](A){\scriptsize$A$}
\tkzLabelPoint[below](B){\scriptsize$B$}
\tkzLabelPoint[above right](C){\scriptsize$C$}
\tkzLabelPoint[above right](D){\scriptsize$D$}
\tkzLabelPoint[above left](E){\scriptsize$E$}
\tkzLabelPoint[above left](F){\scriptsize$F$} 
\tkzLabelPoint[below](G){\scriptsize$G$}
\tkzLabelPoint[below](H){\scriptsize$H$}
\tkzLabelPoint[above ](I){\scriptsize$I$}
\tkzDrawSegments[dashed, line width=1pt](G,H G,I I,H E,B F,A A,D B,C F,C E,D) 
\tkzDrawSegments[->, >=latex,color=red,line width=1.2pt](O,C O,D)
\tkzLabelSegment[below](O,C){$\vv{\imath}$}
\tkzLabelSegment[below](O,D){$\vv{\jmath}$} 
\end{tikzpicture}
}
\hfill
\parbox{0.32\linewidth}{
\begin{tikzpicture} 
\tkzSetUpPoint[size=3, fill=black]
\tkzfctset{tan style/.style={->,>=latex, color=red, line width=1.2pt}}   
\tkzSetUpLine[line width= 1.2pt, add=.2 and .2] 

\tkzDefPoint(0,0){O}
\def \rayon{1.8};  
\tkzDefShiftPoint[O](0:\rayon){C}
\tkzDefShiftPoint[O](60:\rayon){D}
\tkzDefShiftPoint[O](120:\rayon){E}
\tkzDefShiftPoint[O](180:\rayon){F}
\tkzDefShiftPoint[O](240:\rayon){A}
\tkzDefShiftPoint[O](300:\rayon){B}
\tkzDefShiftPoint[B](0:\rayon){H}
\tkzDefShiftPoint[E](60:\rayon){I}
\tkzDefShiftPoint[A](180:\rayon){G}
%	\tkzDrawPolygon(A,B,C,D,E,F,G,H,I)
\tkzLabelPoint[above](O){\scriptsize$O$}
\tkzLabelPoint[below](A){\scriptsize$A$}
\tkzLabelPoint[below](B){\scriptsize$B$}
\tkzLabelPoint[above right](C){\scriptsize$C$}
\tkzLabelPoint[above right](D){\scriptsize$D$}
\tkzLabelPoint[above left](E){\scriptsize$E$}
\tkzLabelPoint[above left](F){\scriptsize$F$} 
\tkzLabelPoint[below](G){\scriptsize$G$}
\tkzLabelPoint[below](H){\scriptsize$H$}
\tkzLabelPoint[above ](I){\scriptsize$I$}
\tkzDrawSegments[dashed, line width=1pt](G,H G,I I,H E,B F,A A,D B,C F,C E,D) 
\tkzDrawSegments[->, >=latex,color=red,line width=1.2pt](O,C O,D)
\tkzLabelSegment[below](O,C){$\vv{\imath}$}
\tkzLabelSegment[below](O,D){$\vv{\jmath}$} 
\end{tikzpicture}
}\hfill
\parbox{0.32\linewidth}{
\begin{tikzpicture} 
\tkzSetUpPoint[size=3, fill=black]
\tkzfctset{tan style/.style={->,>=latex, color=red, line width=1.2pt}}   
\tkzSetUpLine[line width= 1.2pt, add=.2 and .2] 

\tkzDefPoint(0,0){O}
\def \rayon{1.8};  
\tkzDefShiftPoint[O](0:\rayon){C}
\tkzDefShiftPoint[O](60:\rayon){D}
\tkzDefShiftPoint[O](120:\rayon){E}
\tkzDefShiftPoint[O](180:\rayon){F}
\tkzDefShiftPoint[O](240:\rayon){A}
\tkzDefShiftPoint[O](300:\rayon){B}
\tkzDefShiftPoint[B](0:\rayon){H}
\tkzDefShiftPoint[E](60:\rayon){I}
\tkzDefShiftPoint[A](180:\rayon){G}
%	\tkzDrawPolygon(A,B,C,D,E,F,G,H,I)
\tkzLabelPoint[above](O){\scriptsize$O$}
\tkzLabelPoint[below](A){\scriptsize$A$}
\tkzLabelPoint[below](B){\scriptsize$B$}
\tkzLabelPoint[above right](C){\scriptsize$C$}
\tkzLabelPoint[above right](D){\scriptsize$D$}
\tkzLabelPoint[above left](E){\scriptsize$E$}
\tkzLabelPoint[above left](F){\scriptsize$F$} 
\tkzLabelPoint[below](G){\scriptsize$G$}
\tkzLabelPoint[below](H){\scriptsize$H$}
\tkzLabelPoint[above ](I){\scriptsize$I$}
\tkzDrawSegments[dashed, line width=1pt](G,H G,I I,H E,B F,A A,D B,C F,C E,D) 
\tkzDrawSegments[->, >=latex,color=red,line width=1.2pt](O,C O,D)
\tkzLabelSegment[below](O,C){$\vv{\imath}$}
\tkzLabelSegment[below](O,D){$\vv{\jmath}$} 
\end{tikzpicture}
}
\begin{multicols}{3}
\begin{enumerate}[label={\alph*)\rule{0pt}{2.1em}},leftmargin=*] 
\item $\vv{BI}= -2\vv{\imath}+3\vv{\jmath}$
\item $\vv{IA}= \vv{\imath}-3\vv{\jmath}$
\item $\vv{IH}= 3\vv{\imath}-3\vv{\jmath}$
\item $\vv{EB}= 2\vv{\imath}-2\vv{\jmath}$
\item $\vv{OG}= -\vv{\imath}-\vv{\jmath}$
\item $\vv{DG}= -\vv{\imath}-2\vv{\jmath}$
\item $\vv{GC}= 2\vv{\imath}+\vv{\jmath}$
\item $\vv{HE}= -3\vv{\imath}+2\vv{\jmath}$ 
\item $\vv{AE}= -\vv{\imath}+2\vv{\jmath}$ 
\item $\vv{AF}= -\vv{\imath}+\vv{\jmath}$
\item $\vv{BD}= -\vv{\imath}+2\vv{\jmath}$
\item $\vv{FB}= 2\vv{\imath}-\vv{\jmath}$ 
\end{enumerate}
\end{multicols}


\end{proof}

\begin{proof}[solution de \ref{chapter16:vecteur:decompositionB}]\hfill 


\noindent
\parbox{0.32\linewidth}{	
\begin{tikzpicture} 
\tkzSetUpPoint[size=3, fill=black]
\tkzfctset{tan style/.style={->,>=latex, color=red, line width=1.2pt}}   
\tkzSetUpLine[line width= 1.2pt, add=.2 and .2] 

\tkzDefPoint(0,0){O}
\def \rayon{2.3};  
\tkzDefShiftPoint[O](-120:\rayon){A}
\tkzDefShiftPoint[O](-60:\rayon){B}
\tkzDefShiftPoint[O](0:\rayon){C}
\tkzDefShiftPoint[O](60:\rayon){D}
\tkzDefShiftPoint[O](120:\rayon){E}
\tkzDefShiftPoint[O](180:\rayon){F}
\tkzDrawPoints(O,A,B,C,D,E,F)
\tkzDrawPolygon(A,B,C,D,E,F)  
\tkzLabelPoint[below](O){\scriptsize$O$}
\tkzLabelPoint[below](A){\scriptsize$A$}
\tkzLabelPoint[below](B){\scriptsize$B$}
\tkzLabelPoint[right](C){\scriptsize$C$}
\tkzLabelPoint[above](D){\scriptsize$D$}
\tkzLabelPoint[above](E){\scriptsize$E$}
\tkzLabelPoint[left](F){\scriptsize$F$}
\tkzDefShiftPoint[O](0:{0.5*\rayon}){K}
\tkzDefShiftPoint[O](60:{0.5*\rayon}){L}
\tkzDefShiftPoint[O](120:{0.5*\rayon}){M}
\tkzDefShiftPoint[O](180:{0.5*\rayon}){N}
\tkzDefShiftPoint[O](240:{0.5*\rayon}){I}
\tkzDefShiftPoint[O](300:{0.5*\rayon}){J}
\tkzDrawPolygon(I,J,K,L,M,N)
\tkzLabelPoint[below](I){\scriptsize$I$}
\tkzLabelPoint[below](J){\scriptsize$J$}
\tkzLabelPoint[above](M){\scriptsize$M$}
\tkzLabelPoint[above](L){\scriptsize$L$}
\tkzLabelPoint[above right](K){\scriptsize$K$}
\tkzLabelPoint[above left](N){\scriptsize$N$}
\tkzDrawSegments[dashed, line width=1pt](O,A O,B O,C O,D O,E O,F)  
\tkzDrawSegments[->, >=latex, line width=1.2pt](O,I O,J)
%		\tkzLabelSegment[below](O,I){$\vv{\imath}$}
%		\tkzLabelSegment[below](O,J){$\vv{\jmath}$} 
\end{tikzpicture}	
} \hfill 
\parbox{0.65\linewidth}{
Décomposer chacun des vecteurs suivants selon les vecteurs donnés.
\begin{multicols}{2}
\begin{enumerate}[label={\alph*)\rule{0pt}{2.1em}},leftmargin=*] 
\item $\vv{BF}= 2\vv{OI}-4\vv{OJ}$
\item $\vv{JD}=-3\vv{OI}-\vv{OK}$
\item $\vv{CO}=2\vv{OI}-2\vv{OJ}$
\item $\vv{EA}=4\vv{OI}+2\vv{OK}$
\end{enumerate}
\end{multicols}
}


\noindent
\parbox{0.32\linewidth}{	
\begin{tikzpicture} 
\tkzSetUpPoint[size=3, fill=black]
\tkzfctset{tan style/.style={->,>=latex, color=red, line width=1.2pt}}   
\tkzSetUpLine[line width= 1.2pt, add=.2 and .2] 

\tkzDefPoint(0,0){O}
\def \rayon{2.3}; 
\tkzDefShiftPoint[O](-120:\rayon){A}
\tkzDefShiftPoint[O](-60:\rayon){B}
\tkzDefShiftPoint[O](0:\rayon){C}
\tkzDefShiftPoint[O](60:\rayon){D}
\tkzDefShiftPoint[O](120:\rayon){E}
\tkzDefShiftPoint[O](180:\rayon){F}
\tkzDrawPoints(O,A,B,C,D,E,F)
\tkzDrawPolygon(A,B,C,D,E,F)  
\tkzLabelPoint[below](O){\scriptsize$O$}
\tkzLabelPoint[below](A){\scriptsize$A$}
\tkzLabelPoint[below](B){\scriptsize$B$}
\tkzLabelPoint[right](C){\scriptsize$C$}
\tkzLabelPoint[above](D){\scriptsize$D$}
\tkzLabelPoint[above](E){\scriptsize$E$}
\tkzLabelPoint[left](F){\scriptsize$F$}
\tkzDefShiftPoint[O](0:{0.5*\rayon}){K}
\tkzDefShiftPoint[O](60:{0.5*\rayon}){L}
\tkzDefShiftPoint[O](120:{0.5*\rayon}){M}
\tkzDefShiftPoint[O](180:{0.5*\rayon}){N}
\tkzDefShiftPoint[O](240:{0.5*\rayon}){I}
\tkzDefShiftPoint[O](300:{0.5*\rayon}){J}
\tkzDrawPolygon(I,J,K,L,M,N)
\tkzLabelPoint[below](I){\scriptsize$I$}
\tkzLabelPoint[below](J){\scriptsize$J$}
\tkzLabelPoint[above](M){\scriptsize$M$}
\tkzLabelPoint[above](L){\scriptsize$L$}
\tkzLabelPoint[above right](K){\scriptsize$K$}
\tkzLabelPoint[above left](N){\scriptsize$N$}
\tkzDrawSegments[dashed, line width=1pt](O,A O,B O,C O,D O,E O,F)  
\tkzDrawSegments[->, >=latex, line width=1.2pt](O,A O,B)
%		\tkzLabelSegment[below](O,I){$\vv{\imath}$}
%		\tkzLabelSegment[below](O,J){$\vv{\jmath}$} 
\end{tikzpicture}	
} \hfill 
\parbox{0.65\linewidth}{
Écrire chacun des vecteurs suivants sous la forme $a\vv{OA}+b\vv{OB}$.
\begin{multicols}{2}
\begin{enumerate}[label={\alph*)\rule{0pt}{2.1em}},leftmargin=*] 
\item $\vv{JF}= \vv{OA}-\tfrac{3}{2}\vv{OB}$  
\item $\vv{EK}= \vv{OA}+\tfrac{3}{2}\vv{OC}$
\item $\vv{CI}= \tfrac{3}{2}\vv{OA}-\vv{OB}$
\item $\vv{JD}= -\tfrac{3}{2}\vv{OA}-\tfrac{1}{2}\vv{OC}$  
\end{enumerate}
\end{multicols}
} 


\end{proof}

%%----------------------------------------------------------------------------------------
%%	 
%%----------------------------------------------------------------------------------------
%
%\newpage 
%\loadgeometry{marginlayout}  
%\pagestyle{scrheadings} % Print headers again  
%\KOMAoptions{
%headwidth=textwithmarginpar,
%headsepline=true,
%headsepline= 0.5pt,
%footwidth=textwithmarginpar,
%} 
%\onehalfspacing
%
%\section{Produit d'un vecteur par un réel}
%
%\begin{definition}[interprétation géométrique] $\vv{u}\neq \vv{0}$. 
%Le produit d'un réel $k$ par un vecteur $\vv{u}$ est un vecteur noté $k\vv{u}$ : 
%\begin{description}
%\item[P1]  $0\vv{u}=\vv{0}$.
%\item[P2]  Si $k \neq  0$,  \og $k\vv{u}$ \fg{} désigne le vecteur :
%\begin{enumerate}[label=\textbullet, leftmargin=*]
%\item ayant même direction que $\vv{u}$ 
%\item $\norm{k\vv{u}}=\abs{k}\times \norm{\vv{u}}$
%\item Si $k>0$ alors $k\vv{u}$ et $\vv{u}$ ont \textbf{même sens}.\\
%Si $k<0$ alors $k\vv{u}$ et $\vv{u}$ sont de \textbf{sens contraires}
%\end{enumerate}
%\end{description} 
%\end{definition}
%
%
%\vfill
%
%\begin{definition}[vecteurs opposés]  $\vv{v}=-\vv{u}$ signifie que les vecteurs $\vv{u}$ et $\vv{v}$  ont  même direction ,  même norme mais sont de sens contraires. 
%En particulier :
%\end{definition}
%\begin{tBox}
%$-\vv{AB}=\vv{BA}$
%\end{tBox}
%\begin{example} $I$ est le milieu du segment $[AB]$. En utilisant uniquement des vecteurs d'extrémités  $A$, $B$ ou $I$ :
%\begin{enumerate}[label=\alph*), leftmargin=*]
%\item Énumérer 2 égalités de vecteurs.
%\vfill
%\item Énumérer 3 paires de vecteurs opposés $\vv{v}=-\vv{u}$
%\vfill
%\item Énumérer 3 paires de vecteurs qui peuvent s'écrire $\vv{v}=2\vv{u}$
%\vfill
%\item Énumérer 3 paires de vecteurs qui peuvent s'écrire $\vv{v}=\frac{1}{2}\vv{u}$
%\vfill
%\end{enumerate}
%\end{example}
%\vfill 
%\begin{tBox} Soit un repère $(O;I,J)$, et le vecteur $\vv{u}\begin{pmatrix}
%x\\y
%\end{pmatrix}$.
%
%\noindent	Le vecteur $k\vv{u}$ a pour coordonnées $k\vv{u}\begin{pmatrix}
%kx\\ky
%\end{pmatrix}$. 
%%\noindent On pose $\vec{\imath}=\vv{OI}\begin{pmatrix}
%%	1\\0
%%\end{pmatrix}$ et $\vec{\jmath}=\vv{OJ}\begin{pmatrix}
%%0\\1
%%\end{pmatrix}$. 
%%
%%\noindent Le couple de vecteurs $(\vec{\imath}, \vec{\jmath}=\vv{OJ})$ est une base. Le vecteur $\vv{u}\begin{pmatrix}
%%	x\\y
%%\end{pmatrix}$ s'écrit comme la combinaison linéaire  $\vv{u}=x\vec{\imath}+y\vec{\jmath}$.
%\end{tBox}   
%%----------------------------------------------------------------------------------------
%%	 
%%----------------------------------------------------------------------------------------
%
%\newpage 
%\loadgeometry{widelayout}  
%\pagestyle{scrheadings} % Print headers again  
%\KOMAoptions{
%headwidth=textwithmarginpar,
%headsepline=true,
%headsepline= 0.5pt,
%footwidth=textwithmarginpar,
%} 
%\onehalfspacing
%
%
%\subsection*{Exercices : produit par un réel}
%\begin{example}\hfill \\
%Les vecteurs ci-dessous ont tous la même direction.\\
%\begin{tikzpicture}[scale=1]
%\tkzInit[xmax=6,ymax=1.5]%\tkzGrid[sub]
%\tkzSetUpPoint[size=4, fill=black]
%\tkzfctset{tan style/.style={->,>=latex, color=red, line width=1.2pt}}  
%
%\def \xa{3}; 
%\def \ya{0}; 
%
%\def \xb{5};
%\def \yb{0};
%
%\tkzDefPoint(1,0){A} \tkzDefShiftPoint[A](\xa,\ya){B}
%\tkzDefPoint(0,-1){C} \tkzDefShiftPoint[C](\xb,\yb){D}
%
%\tkzDrawSegments[->, >=latex, line width=1.2pt](A,B C,D)
%
%
%\tkzLabelPoints[below, font=\scriptsize](C,D) 
%\tkzLabelPoints[above, font=\scriptsize](A,B) 
%
%\foreach \i in {1,2}{%
%\tkzMarkSegment[pos=\i/3,mark=|,  size=2pt](A,B)
%}
%
%\foreach \i in {1,...,4}{%
%\tkzMarkSegment[pos=\i/5,mark=|,  size=2pt](C,D)
%} 
%\end{tikzpicture} 
%\hfill 
%\begin{tikzpicture}[scale=1]
%\tkzInit[xmax=6,ymax=1.5]%\tkzGrid[sub]
%\tkzSetUpPoint[size=4, fill=black]
%\tkzfctset{tan style/.style={->,>=latex, color=red, line width=1.2pt}}  
%
%\def \xa{3.6}; 
%\def \ya{0}; 
%
%\def \xb{2.4};
%\def \yb{0};
%
%\tkzDefPoint(0,0){A} \tkzDefShiftPoint[A](\xa,\ya){B}
%\tkzDefPoint(3,-1){C} \tkzDefShiftPoint[C](-\xb,\yb){D}
%\tkzDrawSegments[->, >=latex, line width=1.2pt](A,B C,D)
%
%\tkzLabelPoints[below, font=\scriptsize](C,D) 
%\tkzLabelPoints[above, font=\scriptsize](A,B) 
%
%\foreach \i in {1,2}{%
%\tkzMarkSegment[pos=\i/3,mark=|,  size=2pt](A,B)
%}
%
%\foreach \i in {1}{%
%\tkzMarkSegment[pos=\i/2,mark=|,  size=2pt](C,D)
%} 
%\end{tikzpicture} 
%\hfill 
%\begin{tikzpicture}[scale=1]
%\tkzInit[xmax=6,ymax=1.5]%\tkzGrid[sub]
%\tkzSetUpPoint[size=4, fill=black]
%\tkzfctset{tan style/.style={->,>=latex, color=red, line width=1.2pt}}  
%
%\def \xa{3}; 
%\def \ya{0}; 
%
%\def \xb{3};
%\def \yb{0};
%
%\tkzDefPoint(0,0){A} \tkzDefShiftPoint[A](\xa,\ya){B}
%\tkzDefPoint(2,-1){C} \tkzDefShiftPoint[C](-\xb,\yb){D}
%\tkzDrawSegments[->, >=latex, line width=1.2pt](A,B C,D)
%\tkzDrawSegments[dashed, line width=0.5pt](A,D C,B)
%\tkzLabelPoints[below, font=\scriptsize](C,D) 
%\tkzLabelPoints[above, font=\scriptsize](A,B) 
%\end{tikzpicture}  
%\begin{minipage}{0.3\columnwidth} 
%\begin{align*}
%\frac{CD}{AB}&=\frac{5}{3} \\
%CD & = \tfrac{5}{3}AB \\
%\vv{CD}& =\tfrac{5}{3}\vv{AB}   
%\end{align*}
%\end{minipage}
%\end{example}
%
%\begin{exercise}\hfill 
%\begin{enumerate}[label=\alph*),leftmargin=*]
%\item Sur la figure ci-dessous : 
%\begin{minipage}{0.8\columnwidth}%
%\centering
%\begin{tikzpicture} 
%\tkzInit[xmin=-3.5, xmax=6.5]
%\tkzDrawX[label={}, line width=1.2pt] %noticks, 
%\tkzSetUpPoint[size=7, fill=black]
%\tkzfctset{tan style/.style={->,>=latex, color=red, line width=1.2pt}} 
%\foreach \x/\xtext in {-3/A, -1/B, 6/C} {
%\draw (\x,0.1cm) -- (\x,-0.1cm) node[above] {$\xtext\strut$};
%} 
%\end{tikzpicture}
%\end{minipage} 
%%	
%\begin{align*}
%{AC}&= \ldots \ldots  {BC}
%& {BC}&=\ldots \ldots  {AC}
%&  {BC}&=\ldots \ldots  {BA} \\
%\vv{AC}&= \ldots \ldots \vv{BC}
%&\vv{BC}&=\ldots \ldots \vv{AC}
%& \vv{BC}&=\ldots \ldots \vv{BA} \\
%{CA}&=\ldots \ldots {BA}
%& {AB}&=\ldots \ldots {BC} 
%& {BA}&=\ldots \ldots {BC} \\
%\vv{CA}&=\ldots \ldots \vv{BA}
%& \vv{AB}&=\ldots \ldots \vv{BC} 
%& \vv{BA}&=\ldots \ldots \vv{BC}  
%\end{align*}  
%\item À partir de la figure ci-dessous, compléter les égalités vectorielles suivantes:
%\\
%\begin{minipage}{1\columnwidth}%
%\centering
%\begin{tikzpicture} 
%\tkzInit[xmin=-3.5, xmax=6.5]
%\tkzDrawX[label={}, line width=1.2pt] %noticks, 
%\foreach \x/\xtext in {-2/A, 1/B, 6/C} {
%\draw[thick] (\x,0.1cm) -- (\x,-0.1cm) node[above] {$\xtext\strut$};
%}
%\end{tikzpicture}
%\end{minipage} 
%\begin{align*}
%\vv{AC}& = \ldots \ldots \vv{BC}
%&\vv{BC}& = \ldots \ldots \vv{BA}
%& \vv{CA}& = \ldots \ldots \vv{AB}\\
%% 		\\
%\vv{BC}& = \ldots \ldots \vv{AC}
%& \vv{BA}& = \ldots \ldots \vv{BC}
%& \vv{BA}& = \ldots \ldots \vv{AB}
%\end{align*} 
%\item Le point $B$ est le milieu de $[AC]$. Compléter les égalités vectorielles suivantes :
%\begin{align*}
%\vv{BA}& = \ldots \ldots \vv{AC}
%&\vv{CB}& = \ldots \ldots \vv{AC}
%& \vv{AC}& = \ldots \ldots \vv{AB} 
%\end{align*}  
%\end{enumerate}
%\end{exercise}
%
%\begin{exercise}\hfill 
%
%Placer les points $M$, $N$, $P$ pour chacun des cas suivants  :
%\begin{enumerate}[label=\alph*),leftmargin=*]
%\item $\vv{AM}=2\vv{AB}$; $\vv{AN}=-3\vv{AB}$ et $\vv{OP}=4\vv{AB}$
%
%\begin{tikzpicture}[scale=0.8]
%\tkzInit[xmax=6,ymax=1.5]%\tkzGrid[sub]
%\tkzSetUpPoint[size=4, fill=black]
%\tkzfctset{tan style/.style={->,>=latex, color=red, line width=1.2pt}}  
%
%\def \xo{0.5}; 
%\def \xa{4}; 
%\def \xb{6};  
%\tkzDefPoint(0,0){O}
%\tkzDefShiftPoint[O](\xa,0){A} 
%\tkzDefShiftPoint[O](\xb,0){B} 
%\tkzDefPoint(12,0){C}  
%\tkzDefPoint(-6,0){D}  
%\tkzDrawSegments[-, >=latex, line width=1.2pt](C,D)
%
%
%\tkzLabelPoints[above, font=\scriptsize](A,B,O) 
%
%\foreach \i in {0,...,18}{%
%%		\tkzMarkSegment[pos=\i/4,mark=|,  size=3pt](A,B)
%%		\tkzMarkSegment[pos=\i/6,mark=|,  size=3pt](A,D)
%\tkzMarkSegment[pos=\i/18,mark=|,  size=3pt](C,D)
%} 
%\end{tikzpicture}
%
%\vspace{1cm}
%\item $\vv{AM}=\tfrac{2}{5}\vv{AB}$; $\vv{BN}=-\tfrac{7}{5}\vv{AB}$ et $\vv{OP}=-\tfrac{1}{5}\vv{AB}$
%
%\begin{tikzpicture}[scale=0.8]
%\tkzInit[xmax=6,ymax=1.5]%\tkzGrid[sub]
%\tkzSetUpPoint[size=4, fill=black]
%\tkzfctset{tan style/.style={->,>=latex, color=red, line width=1.2pt}}  
%
%\def \xo{2.5}; 
%\def \xa{1}; 
%\def \xb{6};  
%\tkzDefPoint(0,0){O}
%\tkzDefShiftPoint[O](\xa,0){A} 
%\tkzDefShiftPoint[O](\xb,0){B} 
%\tkzDefPoint(12,0){C}  
%\tkzDefPoint(-6,0){D}  
%\tkzDrawSegments[-, >=latex, line width=1.2pt](C,D)
%
%
%\tkzLabelPoints[above, font=\scriptsize](A,B,O) 
%
%\foreach \i in {0,...,18}{%
%%		\tkzMarkSegment[pos=\i/4,mark=|,  size=3pt](A,B)
%%		\tkzMarkSegment[pos=\i/6,mark=|,  size=3pt](A,D)
%\tkzMarkSegment[pos=\i/18,mark=|,  size=3pt](C,D)
%} 
%\end{tikzpicture}	
%
%\vspace{1cm}
%\item $\vv{AM}=\tfrac{5}{3}\vv{AB}$; $\vv{BN}=-\tfrac{1}{6}\vv{AB}$ et $\vv{OP}=-\tfrac{4}{3}\vv{AB}$ 
%
%
%\vspace{1cm}
%\begin{tikzpicture}[scale=0.8]
%\tkzInit[xmax=6,ymax=1.5]%\tkzGrid[sub]
%\tkzSetUpPoint[size=4, fill=black]
%\tkzfctset{tan style/.style={->,>=latex, color=red, line width=1.2pt}}  
%
%\def \xo{5.0}; 
%\def \xa{0}; 
%\def \xb{6};  
%\tkzDefPoint(\xo,0){O}
%\tkzDefPoint(\xa,0){A} 
%\tkzDefPoint(\xb,0){B} 
%\tkzDefPoint(12,0){C}  
%\tkzDefPoint(-6,0){D}  
%\tkzDrawSegments[-, >=latex, line width=1.2pt](C,D)
%
%
%\tkzLabelPoints[above, font=\scriptsize](A,B,O) 
%
%\foreach \i in {0,...,18}{%
%%		\tkzMarkSegment[pos=\i/4,mark=|,  size=3pt](A,B)
%%		\tkzMarkSegment[pos=\i/6,mark=|,  size=3pt](A,D)
%\tkzMarkSegment[pos=\i/18,mark=|,  size=3pt](C,D)
%} 
%\end{tikzpicture}	
%
%\vspace{0.5cm}
%\item $2\vv{AM}=-\vv{AB}$; $4\vv{BN}=3\vv{AB}$ et $2\vv{OP}=-3\vv{AB}$
%
%
%\vspace{0.5cm}
%\begin{tikzpicture}[scale=0.8]
%\tkzInit[xmax=6,ymax=1.5]%\tkzGrid[sub]
%\tkzSetUpPoint[size=4, fill=black]
%\tkzfctset{tan style/.style={->,>=latex, color=red, line width=1.2pt}}  
%
%\def \xo{0.0}; 
%\def \xa{3}; 
%\def \xb{7};  
%\tkzDefPoint(\xo,0){O}
%\tkzDefPoint(\xa,0){A} 
%\tkzDefPoint(\xb,0){B}  
%\tkzDefPoint(12,0){C}  
%\tkzDefPoint(-6,0){D}  
%\tkzDrawSegments[-, >=latex, line width=1.2pt](C,D)
%
%
%\tkzLabelPoints[above, font=\scriptsize](A,B,O) 
%
%\foreach \i in {0,...,18}{%
%%		\tkzMarkSegment[pos=\i/4,mark=|,  size=3pt](A,B)
%%		\tkzMarkSegment[pos=\i/6,mark=|,  size=3pt](A,D)
%\tkzMarkSegment[pos=\i/18,mark=|,  size=3pt](C,D)
%} 
%\end{tikzpicture}	
%
%
%\vspace{0.5cm}
%\item $3\vv{AM}=2\vv{AB}$; $6\vv{BN}=11\vv{BA}$ et $6\vv{OP}=-\vv{BA}$
%
%
%\vspace{0.5cm}
%\begin{tikzpicture}[scale=0.8]
%\tkzInit[xmax=6,ymin=-2,ymax=1.5]%\tkzGrid[sub]
%\tkzSetUpPoint[size=4, fill=black]
%\tkzfctset{tan style/.style={->,>=latex, color=red, line width=1.2pt}}  
%
%\def \xo{5.0}; 
%\def \xa{0}; 
%\def \xb{6};  
%\tkzDefPoint(\xo,0){O}
%\tkzDefPoint(\xa,0){A} 
%\tkzDefPoint(\xb,0){B}  
%\tkzDefPoint(12,0){C}  
%\tkzDefPoint(-6,0){D}  
%\tkzDrawSegments[-, >=latex, line width=1.2pt](C,D)
%
%
%\tkzLabelPoints[above, font=\scriptsize](A,B,O) 
%
%\foreach \i in {0,...,18}{%
%%		\tkzMarkSegment[pos=\i/4,mark=|,  size=3pt](A,B)
%%		\tkzMarkSegment[pos=\i/6,mark=|,  size=3pt](A,D)
%\tkzMarkSegment[pos=\i/18,mark=|,  size=3pt](C,D)
%} 
%\end{tikzpicture}	
%
%
%\vspace{0.5cm}
%
%%		\item $\vv{AP}=\frac{7}{2}\vv{AB}$
%%		\item $\vv{AW}=-\frac{9}{2}\vv{AB}$
%%		\item $\vv{AS}=-\frac{4}{5}\vv{AB}$
%%		\item $\vv{AT}=-\frac{1}{4}\vv{AB}$
%%		\item $\vv{AU}= \frac{12}{5}\vv{AB}$
%%		\item $\vv{AV}= \frac{3}{2}\vv{AB}$
%%		\item $\vv{BM}=\frac{12}{5}\vv{AB}$
%%		\item $\vv{BN}=4\vv{AB}$
%%		\item $\vv{BP}=\frac{8}{5}\vv{AB}$
%%		\item $\vv{BQ}=-3\vv{AB}$
%%		\item $\vv{BS}=-\frac{6}{5}\vv{AB}$
%%		\item $\vv{BT}=-\frac{1}{4}\vv{AB}$ 
%\end{enumerate}
%\end{exercise} 
%%\begin{tikzpicture}[scale=0.8]
%%	\tkzInit[xmax=6,ymax=1.5]%\tkzGrid[sub]
%%	\tkzSetUpPoint[size=4, fill=black]
%%	\tkzfctset{tan style/.style={->,>=latex, color=red, line width=1.2pt}}  
%%	
%%	\def \xa{6}; 
%%	\def \ya{0};  
%%	\tkzDefPoint(0,0){A} \tkzDefShiftPoint[A](\xa,\ya){B} 
%%	\tkzDefPoint(12,0){C}  
%%	\tkzDefPoint(-6,0){D}  
%%	\tkzDrawSegments[-, >=latex, line width=1.2pt](C,D)
%%	
%%	
%%	\tkzLabelPoints[above, font=\scriptsize](A,B) 
%%	
%%	\foreach \i in {0,...,6}{%
%%		\tkzMarkSegment[pos=\i/6,mark=|,  size=3pt](A,B)
%%		\tkzMarkSegment[pos=\i/6,mark=|,  size=3pt](A,D)
%%		\tkzMarkSegment[pos=\i/6,mark=|,  size=3pt](B,C)
%%	} 
%%\end{tikzpicture} 
%
%
%\begin{example}\hfill  
%
%\noindent	\begin{minipage}{1.0\columnwidth}\centering
%\begin{tikzpicture}[scale=0.95]
%
%\def \xmin{-7}; 
%\def \xmax{11}; 
%
%\def \ymin{-3};
%\def \ymax{4};	
%\def \alpha{75}
%
%\pgftransformcm{1}{0}{1*cos(\alpha)}{1*sin(\alpha)}{\pgfpointorigin}
%\path[clip,preaction = {draw=gray}] 
%({\xmax-\ymin*cos(\alpha)} , {\ymin}) -- ({\xmax-\ymax*cos(\alpha)},{\ymax}) -- ({\xmin-\ymax*cos(\alpha)},{\ymax}) -- ({\xmin-\ymin*cos(\alpha)},{\ymin}) -- cycle; 
%\tkzInit[xmin=-8, ymin=-8,  xmax=8.5,ymax=9.5]
%\begin{scope}[dashed] 
%\tkzGrid[ xstep=1, ystep=1](-11,-9)(13,15)
%\end{scope}
%%	\foreach \i in {-10, -9,...,10} {
%%		\draw[ color =gray, dashed] (-11,-\i-10) -- (11,12-\i);
%%		%		\tkzFct[domain = -10:10]{-\x-\i};
%%	}
%
%\def \xa{5}; 
%\def \ya{0}; 
%
%\def \xb{0};
%\def \yb{2};
%
%\def \xc{2};
%\def \yc{3};
%
%\tkzDefPoint(-7,3){U0} \tkzDefShiftPoint[U0](\xa,\ya){U1}
%\tkzDefPoint(-3,-2){V0} \tkzDefShiftPoint[V0](\xb,\yb){V1} 
%\tkzDefPoint(-6,-2){W0} \tkzDefShiftPoint[W0](\xc,\yc){W1}
%\tkzDrawSegments[->,>=latex,line width =2 pt, color = Maroon](U0,U1 V0,V1  W0,W1)
%\tkzLabelSegment[above](U0,U1){$\vv{v}$}
%\tkzLabelSegment[left](V0,V1){$\vv{u}$}
%\tkzLabelSegment[left](W1,W0){$\vv{w}$}
%
%% 	\tkzDrawX[left=5pt, label=, font=\scriptsize, line width=1.2pt ] %, style=dashed
%% 	\tkzDrawY[below =5pt, label=, font=\scriptsize, line width=1.2pt]  %, style=dashed
%
%
%\tkzDefPoint(0,0){A}
%\tkzDefPoint(4,-2){B}  
%\tkzDrawPoints(A,B)
%\tkzLabelPoints[font=\scriptsize](A,B) 
%
%\end{tikzpicture} 
%\end{minipage}  
%\begin{enumerate}[label=\alph*),leftmargin=*]
%\item Placer les points $A_0$, $A_1$ tel que $\vv{AA_0}=\frac{3}{2}\vv{u}$ et $\vv{AA_1}=-\vv{u}$. 
%\item Placer les points $B_0$, $B_1$ tel que $\vv{BB_0}=\frac{3}{5}\vv{v}$ et $\vv{B_0B_1}=-\frac{7}{5}\vv{v}$.
%\item Placer les points $C_0$, $C_1$, $C_2$ et $C_3$ tel que $\vv{AC_0}=\vv{w}$, $\vv{AC_1}=-\frac{1}{2}\vv{w}$,  $\vv{BC_2}=\frac{2}{3}\vv{w}$ et  $\vv{BC_3}=\frac{5}{3}\vv{w}$
%\item  Compléter : $\vv{A_0A_1}=\ldots\ldots\vv{u}$; $\vv{B_0B_1}=\ldots\ldots\vv{v}$ et $\vv{C_2C_3}=\ldots\ldots\vv{w}$.
%\end{enumerate} 
%\end{example}
%\begin{exercise}\hfill 
%
%Pour chacun des cas suivants, dites s'il existe un réel $k$ tel que  $\vv{PQ}=k\vv{AB}$ puis préciser si les droites $(PQ)$ et $(AB)$ sont parallèles.
%\begin{multicols}{2} 
%\begin{enumerate}[label=\alph*),leftmargin=*]
%\item $\vv{PQ}\begin{pmatrix}
%4\\6
%\end{pmatrix}$ et $\vv{AB}\begin{pmatrix}
%-16\\-24
%\end{pmatrix} $
%\item $\vv{PQ}\begin{pmatrix}
%-3\\5
%\end{pmatrix}$ et $\vv{AB}\begin{pmatrix}
%-9\\14
%\end{pmatrix}$
%\item $P(6;20)$, $Q(18;2)$ et $\vv{AB}\begin{pmatrix}
%28\\-42
%\end{pmatrix}$
%\item $P(1\; ;\;-1)$, $Q(4\; ;\;3)$, $A(-1\; ;\;5)$ et $B(7\; ;\;1)$
%\item $P(1\; ;\;4)$, $Q(-3\; ;\;5)$, $A(5\; ;\;7)$ et $B(9\; ;\;6)$ 
%%			\item $P(1\; ;\;-1)$, $Q(-2\; ;\; 3)$, $A(3\; ;\;1)$ et $B(-3\; ;\;9)$ 
%%			\item $P(-1\; ;\;5)$, $Q(3\; ;\; -4)$, $A(1 \; ;\;1)$ et $B(9\; ;\;-17)$ 
%%			\item $P(2\; ;\;3)$, $Q(-1\; ;\; 3)$, $A(-2\; ;\;3)$ et $B(-11\; ;\;3)$ 
%%			\item $P(3\; ;\;-4)$, $Q(3\; ;\; 5)$, $A(3\; ;\;-1)$ et $B(3\; ;\;-28)$  
%%			\item $P(−2;1)$, $Q(1;3)$, $A(4;5)$ et $B(13;-13)$
%\end{enumerate}
%\end{multicols}
%\end{exercise}
%
%%----------------------------------------------------------------------------------------
%%	 
%%----------------------------------------------------------------------------------------
%
%\newpage 
%\loadgeometry{marginlayout}  
%\pagestyle{scrheadings} % Print headers again  
%\KOMAoptions{
%headwidth=textwithmarginpar,
%headsepline=true,
%headsepline= 0.5pt,
%footwidth=textwithmarginpar,
%} 
%\onehalfspacing
%
%\section{Colinéarité de deux vecteurs}
%\begin{definition}$\vv{u}\neq \vv{0}$.  Les vecteurs colinéaires à $\vv{u}$ sont tous les multiples $k\vv{u}$ ou $k\in\mathbb{R}$.
%
%Il s'agit du vecteur nul $\vv{0}$ et de tous les vecteurs de même direction que $\vv{u}$. 
%\end{definition} 
%\begin{tBox} Deux vecteurs $\vv{u}$ et $\vv{v}$  sont \textbf{colinéaires}\fg{} (en abrégé $\vv{u}\propto\vv{v}$) si l'un est un multiple de l'autre. 
%\end{tBox} 
%%Les vecteurs colinéaires n'ont pas forcément la même longueur, ni le même sens.
%%\begin{tBox} Deux vecteurs $\vv{u}$ et $\vv{v}$ \textbf{ne sont pas colinéaires} signifie :
%%\begin{enumerate}[label=\alph*)]
%%\item $\vv{u}\neq \vv{0}$ et $\vv{v}\neq \vv{0}$
%%\item ils ont des directions différentes
%%\end{enumerate}
%%\end{tBox}
%
%\begin{postulatT} trois points $A$, $B$ et $C$ sont alignés si et seulement si deux parmi les vecteurs $\vv{AB}$, $\vv{AC}$ ou $\vv{BC}$ sont colinéaires.
%\end{postulatT}
%\begin{postulatT} les droites $(AB)$ et  $(CD)$ sont parallèles si et seulement si  les vecteurs $\vv{AB}$ et $\vv{CDC}$ sont colinéaires.
%\end{postulatT}
%\begin{definition}[Colinéarité à l'aide des coordonnées]
%Dans un repère quelconque $(O;\vec{\imath},\vec{\jmath})$. Soit deux vecteurs $\vv{u}\begin{pmatrix}
%x\\y
%\end{pmatrix}$ et $\vv{v}\begin{pmatrix}
%x'\\y'
%\end{pmatrix}$.
%
%On appelle \og déterminant de $\vv{u}$ et $\vv{v}$ \fg{} le nombre noté :
%\[
%\det(\vv{u}\; ; \; \vv{v}) = \mdet{x & x' \\ y & y'} = xy'- x'y
%\]
%\end{definition}
%\begin{theorem} Deux vecteurs $\vv{u}$ et $\vv{v}$ sont colinéaires si et seulement si $\det(\vv{u};\vv{v})=0$.
%\end{theorem}
%\begin{proof}\hfill \vfill 
%\end{proof}
%
%
%%----------------------------------------------------------------------------------------
%%	 
%%----------------------------------------------------------------------------------------
%
%\newpage 
%\loadgeometry{widelayout}  
%\pagestyle{scrheadings} % Print headers again  
%\KOMAoptions{
%headwidth=textwithmarginpar,
%headsepline=true,
%headsepline= 0.5pt,
%footwidth=textwithmarginpar,
%} 
%\onehalfspacing
%
%\subsection*{Exercices Colinéarité}
%\begin{example}
%Calculer les déteriminant des vecteurs suivants :
%\begin{enumerate}[label=\alph*),leftmargin=*]
%\item \rule{0pt}{2em}
%$\vv{u}\begin{pmatrix}
%3\\2
%\end{pmatrix}$ et $\vv{v}\begin{pmatrix}
%4\\5
%\end{pmatrix}$. \quad $\det(\vv{u};\vv{v})=\mdet{3 & 4 \\ 2 & 5}=$
%\item  \rule{0pt}{2em}
%$\vv{u}\begin{pmatrix}
%-3\\2
%\end{pmatrix}$ et $\vv{v}\begin{pmatrix}
%-1\\4
%\end{pmatrix}$.\quad $\det(\vv{u};\vv{v})=$
%\item  \rule{0pt}{2em}
%$\vv{u}\begin{pmatrix}
%-4\\5
%\end{pmatrix}$ et $\vv{v}\begin{pmatrix}
%8\\-10
%\end{pmatrix}$. 
%\end{enumerate}
%\end{example}
%
%\begin{example}[Interprétation géométrique du déterminant dans un repère orthonormé]\hfill 
%
%Soit un repère \textbf{orthonormé} $(O;I,J)$, et les points $M(a,b)$ et $N(c,d)$ et $P$ tel que $OPMN$ est un parallélogramme. Pour simplifier on suppose que $a$, $b$, $c$ et $d$  sont des réels strictement positifs.
%
%\noindent\parbox{0.5\linewidth}{
%\begin{tikzpicture}[scale=0.8]
%\tkzSetUpPoint[size=3, fill=black]
%\tkzfctset{tan style/.style={->,>=latex, color=red, line width=1.2pt}}   
%\tkzSetUpLine[line width= 1.2pt, add=.2 and .2] 
%\tkzInit[xmin=-0.5,xmax=7.5,ymin=-0.5,ymax=6.5  ]
%\tkzDrawX[  noticks , label=$x$,  font=\small, line width=1.2pt ] %, style=dashed
%\tkzDrawY[  noticks , label=$y$,  font=\small, line width=1.2pt]  %, style=dashed
%\tkzRep[color  = Maroon,line width = 2pt, posxlabel=below, posylabel=left, xnorm=1,ynorm=1, xlabel={}, ylabel={}]
%
%\tkzDefPoint(0,0){O};
%\tkzDefPoint(2,4){N};
%\tkzDefPoint(2,6){xN};
%\tkzDefPoint(0,4){yN};
%\tkzDefPoint(5,2){M};
%\tkzDefPoint(5,0){xM};
%\tkzDefPoint(7,2){yM};
%\tkzDefPoint(7,6){P};
%\tkzDefPoint(7,0){xP};
%\tkzDefPoint(0,6){yP};
%\tkzDrawPolygon[fill=gray!20](O,M,P,N);
%\tkzDrawSegments[dashed](N,xN);
%\tkzDrawSegments[dashed](M,yM);
%%	\tkzDrawSegments[dashed](P,xP P,yP);
%\tkzPointShowCoord(P)
%\tkzPointShowCoord[noydraw](M)
%\tkzPointShowCoord[noxdraw](N)
%\tkzDrawSegment[thick,->, >=latex](O,M);
%\tkzDrawSegment[->,  >=latex  ](O,N);
%\tkzDrawSegment[- , dashed, >=latex](M,P);
%\tkzLabelPoint[above left,font=\small](M){$M(a,b)$}
%\tkzLabelPoint[below right,font=\small](N){$N(c,d)$}
%\tkzLabelPoint[above,font=\small](P){$P(\qquad,\qquad)$}
%\tkzLabelPoints[below left,font=\small](O)
%\end{tikzpicture} 
%}
%
%\end{example}
%\begin{example} Déterminer si les vecteurs $\vv{u}$ et $\vv{v}$ sont colinéaires. 
%\begin{enumerate}[label=\alph*),leftmargin=*]
%\item \rule{0pt}{2em}
%$\vv{u}\begin{pmatrix}
%2\\3
%\end{pmatrix}$ et $\vv{v}\begin{pmatrix}
%6\\9
%\end{pmatrix}$. 
%\item \rule{0pt}{2em}
%$\vv{u}\begin{pmatrix}
%\sqrt{3}+\sqrt{2}\\1
%\end{pmatrix}$ et $\vv{v}\begin{pmatrix}
%1\\\sqrt{3}-\sqrt{2}
%\end{pmatrix}$. 
%\end{enumerate}
%\end{example}
%
%\begin{example}
%Montrer que les points $A(2~;~5)$, $B(3~;~8)$ et $C(-5~;~-16)$ sont alignés.
%\end{example}
%\vfill 
%\begin{example}
%Soient un repère orthonormé et les points $A(0~;~-2)$, $B(1~;~1)$ et $C(2~;~-1)$. Calculer l'aire du triangle $ABC$.
%\end{example}
%\vfill 
%\begin{example}
%Soient $A(1~;~3)$, $B(5~;~-2)$, $C(-1~;~6)$ et $D(7~;~-4)$. 
%Montrer que les droites (AB) et (CD) sont parallèles.
%\end{example}
%\vfill 
%
%\newpage 
%\begin{exercise}
%
%
%Dans chaque cas, dire si les vecteurs sont colinéaires. 
%
%\begin{enumerate}[label=\alph*),leftmargin=*]
%\item  $\vv{u}\begin{pmatrix}
%2\\-3 
%\end{pmatrix}$ et $\vv{v}\begin{pmatrix}
%-1\\\tfrac{3}{2} 
%\end{pmatrix}$  
%\vspace{0.5cm}
%\item  $\vv{u}\begin{pmatrix}
%\tfrac{1}{2}\\\tfrac{1}{3} 
%\end{pmatrix}$ et $\vv{v}\begin{pmatrix}
%\tfrac{4}{5}\\ -\tfrac{1}{3}
%\end{pmatrix}$  
%\vspace{0.5cm}
%\item  $\vv{u}\begin{pmatrix}
%\sqrt{2}\\\sqrt{3} 
%\end{pmatrix}$ et $\vv{v}\begin{pmatrix}
%2\\ -\sqrt{6}
%\end{pmatrix}$  
%\vspace{0.5cm} 
%\end{enumerate}
%\end{exercise}
%\begin{exercise}Dans chaque cas, déterminer le réel $m$ pour que les vecteurs $\vv{u}$ et $\vv{v}$ soient colinéaires.  
%\begin{enumerate}[label=\alph*),leftmargin=*]
%\item  $\vv{u}\begin{pmatrix}
%2\\6
%\end{pmatrix}$ et $\vv{v}\begin{pmatrix}
%m\\ 3
%\end{pmatrix}$  
%\vspace{0.5cm}
%\item  $\vv{u}\begin{pmatrix}
%-m\\4m-3
%\end{pmatrix}$ et $\vv{v}\begin{pmatrix}
%1\\ -3
%\end{pmatrix}$  
%\vspace{0.5cm}
%\item  $\vv{u}\begin{pmatrix}
%27\\2m
%\end{pmatrix}$ et $\vv{v}\begin{pmatrix}
%2m\\ 3
%\end{pmatrix}$  
%\vspace{0.5cm} 
%\end{enumerate} 
%\end{exercise}
%\begin{exercise} 
%Dans un repère orthonormé $( O ; \vec{\imath} , \vec{\jmath} )$.
%\begin{enumerate}[label=\alph*),leftmargin=*]
%\item Montrer que $A\left(-2;1\right)$, $B\left(3;3\right)$ , $C\left(1;\tfrac{11}{5}\right)$ sont alignés.\vfill
%\item Trouver $a$ pour que $A(1 ; 2)$, $B(-3 ; 0)$ et $C(-14 ; a)$ soient alignés.\vfill 
%\end{enumerate} 
%\end{exercise}  
%\begin{exercise}
%
%Dans un repère, soit les points: $A(-2~;~1)$, $B(3~;~3)$, $C\left(1~;~\tfrac{11}{5}\right)$ et $D\left(\tfrac{45}{2}~;~\tfrac{54}{5}\right)$
%\begin{enumerate}[label=\alph*),leftmargin=*]
%\item Démontrer que les points $A$, $B$ et $C$ sont alignés. 
%\vspace{2cm}
%\item Les points $A$, $B$ et $D$ sont-ils alignés ?
%\vspace{2cm}
%\end{enumerate}
%
%\end{exercise}
%\newpage 
%\begin{exercise} Soit les points  $A(1 ; 2)$, $B(-3 ; 0)$ et $C(-7 ; a)$ dans un repère orthonormé. Trouver 2 valeurs de $a$ pour lesquelles le triangle $ABC$ est d'aire $12$.  
%\end{exercise}
%\begin{exercise} \hfill 
%
%Dans un repère, on donne les points : $M(0~;~-3)$, $N(2~;~3)$, $P(-9~;~0)$ et $Q(-1~;~-1)$
%\begin{enumerate}[label=\arabic*),leftmargin=*]
%\item Calculer les coordonnées des points $A$ et $B$ tels que: 
%\begin{enumerate}[label=\alph*),leftmargin=*]
%\item $\vv{NA}=\dfrac{1}{2}~\vv{MN}$
%\vspace{2cm}
%\item $\vv{MB}=3~\vv{MQ}$
%\vspace{2cm}
%\end{enumerate}
%\item Calculer les coordonnées des vecteurs $\vv{PA}$ et $\vv{PB}$.
%\vspace{2cm} 
%\item Démontrer que les points $P$, $A$ et $B$ sont alignés.
%\vspace{3cm}
%\end{enumerate}
%\end{exercise}
%\vfill 






%----------------------------------------------------------------------------------------
%	 
%----------------------------------------------------------------------------------------

\newpage 
\loadgeometry{marginlayout}  
\pagestyle{scrheadings} % Print headers again  
\KOMAoptions{
headwidth=textwithmarginpar,
headsepline=true,
headsepline= 0.5pt,
footwidth=textwithmarginpar,
} 
\onehalfspacing


\section{Propriétés des opérations}   

\begin{tBox} 
\textbf{Élément nul}\quad	Pour tout vecteur $\vv{u}$, \qquad $\vv{0}+\vv{u}=\vv{u}+\vv{0}=\vv{u}$. 
\end{tBox}
\begin{tBox}[frametitle={La somme de 2 vecteurs est indépendante de l'ordre}]
Pour tout vecteur $\vv{u}$ et  $\vv{v}$ \qquad $\vv{u}+\vv{v}=\vv{v}+\vv{u}$.
\end{tBox}

\begin{minipage}{0.47\columnwidth} 
\begin{tikzpicture}[scale=0.85]
%\tkzInit[xmax=6,ymax=1.5]%\tkzGrid[sub]
\tkzSetUpPoint[size=4, fill=black]
\tkzfctset{tan style/.style={->,>=latex, color=red, line width=1.2pt}}  

\def \xa{0.8}; 
\def \ya{2.4}; 

\def \xb{4};
\def \yb{0};

\tkzDefPoint(-0.7,0){U0} \tkzDefShiftPoint[U0](\xa,\ya){U1}
\tkzDefPoint(1.0,3.1){V0} \tkzDefShiftPoint[V0](\xb,\yb){V1}

\tkzDrawSegments[->, >=latex, line width=1.2pt](U0,U1 V0,V1)
\tkzLabelSegment(U0,U1){\small $\vv{u}$}
\tkzLabelSegment(V0,V1){\small $\vv{v}$}

\tkzDefPoint(0.8,0){A}  
\tkzDefPointBy[translation=from U0 to U1](A)\tkzGetPoint{B}
\tkzDefPointBy[translation=from V0 to V1](B)\tkzGetPoint{C}

\tkzDrawSegments[->, >=latex, line width=0.6pt, dashed](A,B)
\tkzDrawSegments[-, line width=0.6pt, dashed](B,C)
\tkzDrawSegments[->, >=latex, line width=1.2pt](A,C)

\tkzLabelPoint[below](A){\scriptsize$A$}
\tkzLabelPoint[above left](B){\scriptsize$B$}
\tkzLabelPoint[right](C){\scriptsize$C$}
\tkzLabelSegment[left](A,B){\small$\vv{AB}$}
\tkzLabelSegment[above](B,C){\small$\vv{BC}$}
\tkzLabelSegment[below, sloped](A,C){\small$\vv{u}+\vv{v}$}


\tkzDefPointBy[translation=from V0 to V1](A)\tkzGetPoint{D} 

\tkzDrawSegments[->, >=latex, line width=0.6pt, dashed](A,D) 
\tkzDrawSegments[line width=0.6pt, dashed](D,C)

\tkzLabelPoint[below right](D){\scriptsize$D$}
\tkzLabelSegment[below](A,D){\small$\vv{AD}$}
\tkzLabelSegment[right](D,C){\small$\vv{DC}$} 
\end{tikzpicture}
\end{minipage} 
\hfill
\begin{minipage}{0.47\columnwidth} 
\begin{align*}
\vv{u}+\vv{v}& = \vv{AB}+\vv{BC} =\vv{AC}\\
\vv{v}+\vv{u}& = \vv{AD}+\vv{DC} = \vv{AC}
\end{align*}
\end{minipage} 
\begin{tBox} 
La somme de plusieurs vecteurs quelconques $\vv{u}$, $\vv{v}$ et $\vv{w}$ est aussi indépendante de l'ordre : 
\end{tBox}



\begin{figure*}[!h] 
\begin{tikzpicture}[scale=1]
\tkzInit[xmax=6,ymax=1.5]%\tkzGrid[sub]
\tkzSetUpPoint[size=4, fill=black]
\tkzfctset{tan style/.style={->,>=latex, color=red, line width=1.2pt}}  

\def \xa{1.0};
\def \ya{1.2};


\def \xb{1.8};
\def \yb{-2.5};

\def \xc{2.2}; 
\def \yc{0.5}; 


\tkzDefPoint(-1.0,0.8){U0} \tkzDefShiftPoint[U0](\xa,\ya){U1}
\tkzDefPoint(0.0,3.5){V0} \tkzDefShiftPoint[V0](\xb,\yb){V1}
\tkzDefPoint(1.0,3.5){W0} \tkzDefShiftPoint[W0](\xc,\yc){W1}

\tkzDrawSegments[->, >=latex, line width=1.2pt](U0,U1 V0,V1 W0,W1)
\tkzLabelSegment(U0,U1){\small $\vv{u}$}
\tkzLabelSegment(V0,V1){\small $\vv{v}$}
\tkzLabelSegment(W0,W1){\small $\vv{w}$}

\tkzDefPoint(4.0,2){A}  
\tkzDefPointBy[translation=from U0 to U1](A)\tkzGetPoint{B}
\tkzDefPointBy[translation=from V0 to V1](B)\tkzGetPoint{C}
\tkzDefPointBy[translation=from W0 to W1](C)\tkzGetPoint{D}


\tkzDrawSegments[->, >=latex, line width=1.0pt](A,C)
\tkzDrawSegments[->, >=latex, line width=2pt](A,D)
\tkzDrawSegments[ line width=0.5pt, dashed](A,B B,C C,D)
\tkzLabelSegment[below , sloped](A,C){\small $\vv{u}+\vv{v}$}
\tkzLabelSegment[below , sloped](C,D){\small $\vv{w}$}

\tkzDefPoint(10,2){E}  
\tkzDefPointBy[translation=from U0 to U1](E)\tkzGetPoint{F}
\tkzDefPointBy[translation=from V0 to V1](F)\tkzGetPoint{G}
\tkzDefPointBy[translation=from W0 to W1](G)\tkzGetPoint{H}



\tkzDrawSegments[->, >=latex, line width=2pt](E,H)
\tkzDrawSegments[->, >=latex, line width=1pt](F,H)
\tkzDrawSegments[ line width=0.5pt, dashed](E,F F,G G,H)

\tkzLabelSegment[above, sloped](E,F){\small $\vv{u}$}
\tkzLabelSegment[above , sloped](F,H){\small $\vv{v}+\vv{w}$}

\tkzText(9,4){$\vv{u}+\vv{v}+\vv{w} = (\vv{u}+\vv{v})+\vv{w} = \vv{u}+(\vv{v}+\vv{w})$}
\end{tikzpicture}
\end{figure*}

\begin{tBox}
Pour tout vecteurs $\vv{u}$,  et réel $a$ et $b$,   \qquad  $a(b\vv{u})=(ab)\vv{u} $. 
\end{tBox}
\begin{tBox}
\vfill 
Pour tout $a$ et $b\in\mathbb{R}$ et vecteur $\vv{u}$ on a : $a\vv{u}+b\vv{u}= (a+b)\vv{u}$
\end{tBox}
\begin{tBox}
Pour tout vecteurs $\vv{u}$, $\vv{v}$  et réel $k$,   \qquad  $k(\vv{u}+\vv{v}  ) =k\vv{u}+k\vv{v} $.
%\begin{align*}
%k(\vv{u}+\vv{v}  )&=k\vv{u}+k\vv{v} 
%%&k(\vv{u}-\vv{v})&=k\vv{u}-k\vv{v}
%\end{align*} 
\end{tBox}
\begin{proof} La figure correspond au cas $k>0$.  Le cas $k<0$ est similaire.\\
\begin{minipage}{0.42\columnwidth} 
\begin{tikzpicture}[scale=0.90]
\tkzInit[xmax=6,ymax=1.5]%\tkzGrid[sub]
\tkzSetUpPoint[size=3, fill=black]
\tkzfctset{tan style/.style={->,>=latex, color=red, line width=1.2pt}}  

\def \xa{-0.22};
\def \ya{1.4};


\def \xb{2.08};
\def \yb{1.04};

\tkzDefPoint(1,0){O}  

\tkzDefPoint(-4.0,-0.5){A} \tkzDefShiftPoint[A](\xa,\ya){B}
\tkzDefShiftPoint[B](\xb,\yb){C} 

\tkzDrawSegments[->, >=latex, line width=1.2pt](A,B B,C)
\tkzDrawSegments[->, style={->,shorten >=3pt,>=latex'}, >=latex, line width=1.2pt](A,C) 
\tkzLabelSegment(A,B){\small $\vv{u}$}
\tkzLabelSegment(B,C){\small $\vv{v}$} 
\tkzLabelSegment[below, sloped](A,C){\small $\vv{u}+\vv{v}$} 


\tkzDefPointBy[homothety=center O ratio 1.5](A) \tkzGetPoint{D}
\tkzDefPointBy[homothety=center O ratio 1.5](B) \tkzGetPoint{E}
\tkzDefPointBy[homothety=center O ratio 1.5](C) \tkzGetPoint{F}
% 
\tkzDrawSegments[->, >=latex, line width=1.2pt](D,E E,F) 
\tkzDrawSegments[->, style={->,shorten >=3pt,>=latex'}, >=latex, line width=1.2pt](D,F) 
\tkzDrawSegments[-, >=latex, line width=0.5pt, dashed](O,D O,E O,F) 
% 
\tkzLabelPoints[below , font=\scriptsize](A,D)
\tkzLabelPoints[above right , font=\scriptsize](C,F)
\tkzLabelPoints[above left , font=\scriptsize](E)
\tkzLabelPoints[above , font=\scriptsize](B)
\tkzLabelPoints[ font=\scriptsize](O)
\tkzDrawPoints(O) 
\end{tikzpicture}
\end{minipage}  
\end{proof}



%----------------------------------------------------------------------------------------
%	 
%----------------------------------------------------------------------------------------

\newpage 
\loadgeometry{widelayout}  
\pagestyle{scrheadings} % Print headers again  
\KOMAoptions{
headwidth=textwithmarginpar,
headsepline=true,
headsepline= 0.5pt,
footwidth=textwithmarginpar,
} 
\onehalfspacing

\subsection*{Exercices bilan opérations sur vecteurs}
\begin{example} \hfill 

\noindent Dans le repère 	$(O~;~ \vec{\imath}~,~ \vec{\jmath})$ on   donne  $\vv{u}\begin{pmatrix}
-2\\5
\end{pmatrix}$ ; $A(1~;~-3)$ et $B(4~;~-5)$.  \\
\noindent Calculer les coordonnées de chacun des vecteurs suivants :\\
\noindent \parbox{0.24\linewidth}{
$-3~\vv{AB}$\\
%\vspace{4cm}	
}
\parbox{0.24\linewidth}{
$\vv{u}+\vv{AB}$\\
%\vspace{4cm}
}
\parbox{0.24\linewidth}{
$-3\vv{u}+2\vv{AB}$ \\
%\vspace{4cm}
}
\parbox{0.24\linewidth}{
$2~\vv{u}+5~\vv{BA}$\\
%	\vspace{4cm}
} 
\end{example} 
\vfill 
%\begin{exercise}\hfill 
%	
%	\noindent Soient $\vv{u}\begin{pmatrix}
%		-2\\5
%	\end{pmatrix}$ et $\vv{v}\begin{pmatrix}
%		3\\ \tfrac{1}{4}
%	\end{pmatrix}$ deux vecteurs définis par leurs coordonnées dans le repère
%	$(O~;~ \vec{\imath}~,~ \vec{\jmath})$.
%	
%	\noindent  Calculer les coordonnées de chacun des vecteurs suivants :
%	\begin{multicols}{5}
%		\begin{enumerate}[label=\alph*),leftmargin=*]
%			\item $\vv{u}+\vv{v}$
%			%\vspace{0.5cm}
%			\item $2\vv{u}$
%			%\vspace{0.5cm}
%			\item $-\dfrac{3}{7}~\vv{v}$
%			%\vspace{0.5cm}
%			\item $-3\vv{u}+\vv{v}$
%			%\vspace{0.5cm}
%			\item $\dfrac{3}{2}~\vv{u}-\dfrac{3}{5}~\vv{v}$
%		\end{enumerate}
%	\end{multicols}
%\end{exercise}
\begin{exercise}\hfill 

\noindent Soient $\vv{u}\begin{pmatrix}
2\\6
\end{pmatrix}$ et $\vv{v}\begin{pmatrix}
-5\\ 3
\end{pmatrix}$ deux vecteurs définis par leurs coordonnées dans le repère
$(O~;~ \vec{\imath}~,~ \vec{\jmath})$.

\noindent  Calculer les coordonnées de chacun des vecteurs suivants :
\begin{multicols}{5}
\begin{enumerate}[label=\alph*),leftmargin=*]
\item $-\vv{u}$
\item $\vv{u}+\vv{v}$
%\vspace{0.5cm}
\item $3\vv{u}+2\vv{v}$
%\vspace{0.5cm}
\item  $-5\vv{u}-\vv{v}$
%\vspace{0.5cm}
%			\item $5\vv{u}-2\vv{v}$
%\vspace{0.5cm}
\item $\tfrac{1}{2}~(\vv{u}- \vv{v})$
\end{enumerate}
\end{multicols}
\end{exercise}
\begin{exercise} \hfill 

\noindent Le plan est muni du repère orthonormé $(O~;~ \vec{\imath}~,~ \vec{\jmath})$ et des points $A(-3~;~2)$, $B(1~;~-2)$, $C(-5~;~3)$.  Déterminer les coordonnées des vecteurs :
\begin{multicols}{5}
\begin{enumerate}[label=\alph*),leftmargin=*]
\item $\vv{AB}$
\item $\vv{AC}$
%\vspace{0.5cm}
\item $\vv{AB}+2\vv{AC}$
%\vspace{0.5cm}
\item  $2\vv{AB}-3\vv{BC}$
%\vspace{0.5cm}
\item $\tfrac{1}{3}\vv{AC}+\tfrac{1}{5}\vv{BC}$ 
\end{enumerate}
\end{multicols}  
\end{exercise}
\begin{exercise}[révision] \hfill 

\noindent Le plan est muni du repère   $(O~;~ \vec{\imath}~,~ \vec{\jmath})$. Soient les points $T(1~;~-2)$, $R(0~;~-1)$, $U(5~;~3)$. Déterminer les coordonnées du point $E(x;y)$ tel que $TRUE$ est un parallélogramme, ainsi que celles du centre $F$ du parallélogramme.
\end{exercise}
\begin{exercise}  Le plan est muni du repère   $(O~;~ \vec{\imath}~,~ \vec{\jmath})$. Soient les points $A(2~;~-3)$, $B(1~;~-1)$, $C(4~;~5)$

\begin{enumerate}[label=\alph*),leftmargin=*]
\item  Le point $M(x;y)$ est tel que $\vv{AM}=-2\vv{BC}$. \\
\noindent Écrire les équations vérifiées par $x$ et $y$ et retrouver les coordonnées de $M$.
%\vspace{3cm}
\item Déterminer les coordonnées du point $N(x;y)$ tel que $\vv{CN}=2~\vv{AB}-3~\vv{AC}$
%\vspace{3cm}
\item Déterminer les coordonnées du point $G(x;y)$ tel que $\vv{GA}+\vv{GB}+\vv{GC}=\vv{0}$
\end{enumerate} 
\end{exercise}





%\begin{exercise}\hfill 
%
%Soient $A(-2;-1)$, $B(3;2)$ et $C(1;5)$ trois points dans un repère orthonormé $(O~;~ \vec{\imath}~,~ \vec{\jmath})$.
%\begin{enumerate}[label=\alph*),leftmargin=*]
%%\item Déterminer les coordonnées du vecteur $\vv{AB}$ puis celles du vecteur $\vv{AC}$.
%\item Déterminer les coordonnées du point $I$ milieu de $[AB]$.
%%\item Calculer la longueur $AC$.
%%\item Déterminer les coordonnées de $\vv{AB}+\vv{AC}$, de $3~\vv{AB}$ et de $2~\vv{AB}+3~\vv{AC}$ .
%\item Déterminer les coordonnées du point $K(x;y)$  tel que  $\vv{CK}=3~\vv{AB}$. 
%\item Déterminer les coordonnées du point $L$ tel que $\vv{AL}+2~\vv{BL}+\vv{CL}=\vv{0}$.
%\end{enumerate} 
%\end{exercise}

\begin{exercise}\hfill 

\noindent Le plan est muni du repère orthonormé $(O~;~ \vec{\imath}~,~ \vec{\jmath})$. Soient les points $A(3~;~-3)$, $B(2~;~-4)$, $C(-5~;~1)$

\begin{enumerate}[label=\alph*),leftmargin=*]
\item Déterminer les coordonnées des vecteurs $\vv{AB}$, $\vv{AC}$ et $\vv{BC}$.
%\vspace{3cm}
\item  Déterminer les coordonnées des vecteurs $\vv{u}=\vv{AC}+2~\vv{AB}$, $\vv{v}=-3\vv{AB}+4\vv{AC}$. 
%\vspace{3cm}
\item Déterminer les coordonnées du point $D$ tel que $\vv{BD}=2~\vv{u}-\vv{v}$.
%\vspace{3cm}
\item Déterminer les coordonnées du point $E$ tel que $\vv{AE}=2~\vv{BE}+2~\vv{u}$
%\vspace{3cm}
\end{enumerate} 
\end{exercise}
\newpage
\begin{exercise}[révision] \hfill 

\noindent	 Soit les points $A(-1~;~3)$; $B(-1~;~2)$; $C(\tfrac{5}{2}~;~2)$ et  $D(-8~;~4)$ dans le repère $(O~;~ \vec{\imath}~,~ \vec{\jmath})$. Montrer que les droites $(AC)$ et $(BD)$ sont parallèles.
\end{exercise}
\begin{exercise}[révision]  \hfill 

\noindent	Soit les points $M(1~;~-2)$; $N(0~;~-1)$ et $P(3~;~-4)$  dans le repère $(O~;~ \vec{\imath}~,~ \vec{\jmath})$. Montrer que les points $M$, $N$ et $P$ sont alignés.
\end{exercise}
\begin{exercise}[révision]  \hfill 

\noindent	Soit les points $A(1~;~1)$; $B(2~;~-3)$; $C(4~;~x)$  dans le repère $(O~;~ \vec{\imath}~,~ \vec{\jmath})$. Trouvez $x$ tel que les points  $A$, $B$ et $C$ sont alignés.
\end{exercise}
\begin{exercise} \hfill 

\noindent 	Dans un repère, on donne les points: $A(1~;~-1)$, $B(-1~;~-2)$ et $C(-2~;~2)$
\begin{enumerate}[label=\arabic*),leftmargin=*]
\item Déterminer les coordonnées du point $G$ vérifiant $\vv{GA}+2\vv{GB}+\vv{GC}=\vv{0}$.
%		\vspace{2.5cm} 
\item Déterminer les coordonnées du point $D$ vérifiant $\vv{BD}=\vv{BA}+\vv{BC}$.
%		\vspace{2.5cm} 
\item Les points $B$, $G$ et $D$ sont-ils alignés ?
\end{enumerate} 
\end{exercise}
\begin{exercise} \hfill 

\noindent 	Dans un repère, on donne les points: $A(-2~;~2)$, $B(0~;~-3)$ et $C(4~;~5)$.
\begin{enumerate}[label=\arabic*),leftmargin=*]
\item Déterminer les coordonnées du point $M$ vérifiant $\vv{AM}=2\vv{AB}-3\vv{AC}$.
%		\vspace{2.5cm} 
\item $I$ est le milieu de $[AB]$. Calculer les coordonnées de $I$
\item Montrer que les points $C$, $I$ et $M$ sont alignés.
\end{enumerate} 
\end{exercise}
\begin{exercise} \hfill 

\noindent 	Dans un repère, on donne les points: $A(2~;~4)$, $B(-2~;~2)$ et $C(6~;~-1)$. $I$ est le milieu de $[AC]$. $G$ et $H$ sont tels que $\vv{AG}=2\vv{AB}-\tfrac{1}{2}\vv{AC}$ et $\vv{BH}=-\tfrac{1}{3}\vv{BC}$.
\begin{enumerate}[label=\arabic*),leftmargin=*]
\item Déterminer les coordonnées de $I$, $G$ et $H$.
\item Prouver que $B$ est le milieu de $[GI]$
\item Montrer que les points $A$, $G$ et $H$ sont alignés.
\end{enumerate} 
\end{exercise}

\begin{exercise}[un classique : approche vectorielle]\hfill 

\noindent Sur la figure ci-contre, $ABCD$ est un carré et les triangles $BCE$ et $ABF$ sont équilatéraux. On considère le repère orthonormé $(A~;~B~,~D)$.\\
\noindent\parbox{0.55\linewidth}{ 
\begin{enumerate}[label=\arabic*),leftmargin=*]
\item Représenter les axes du repère.
\item Donner les coordonnées des points $A$, $B$, $C$ et $D$.
\item Déterminer les coordonnées de $F$ et $E$.\\
\noindent  \emph{indice : trigonométrie} 
\item Démontrer que les points $D$, $F$ et $E$ sont alignés.
\end{enumerate}   
}
\parbox{0.4\linewidth}{
\begin{tikzpicture}
\tkzDefPoints{0/0/A,4/0/B,4/4/C,0/4/D}
\tkzDrawPolygon[line width = 1pt](A,B,C,D)
\tkzDefEquilateral(C,B) \tkzGetPoint{E}
\tkzDefEquilateral(A,B) \tkzGetPoint{F}
\tkzDrawPoints(A,B,C,D,E,F)
\tkzDrawPolygon[line width = 1pt](B,C,E)
\tkzDrawPolygon[line width = 1pt](A,B,F)
\tkzLabelPoints(B)
\tkzLabelPoints[below left](A)
\tkzLabelPoints[above right](C)	
\tkzLabelPoints[above left](D)			
\tkzLabelPoint[right](E){$E$}
\tkzLabelPoint[above](F){$F$}
%\tkzDrawPolygon[dashed](A,E,F)
\end{tikzpicture}
}
\end{exercise}
%%%%%%%%%%%%%%%%%%%%%%%%%%%%%%%%%%%%%%%%%%%%%%%%%%%%%%%%%%%
%
%%%%%%%%%%%%%%%%%%%%%%%%%%%%%%%%%%%%%%%%%%%%%%%%%%%%%%%%%%%
\newpage 
\begin{example}[Équation vectorielle sur une droite]\hfill \\
Soit les points $A$, $B$. Soit $P$ et $Q$  tel que $\vv{PA}+3\vv{PB}=\vv{0}$ et $\vv{QA}-3\vv{QB}=\vv{0}$. Placer $P$ et $Q$.

%\begin{minipage}{0.5\columnwidth}
%%	\begin{enumerate}[label=\alph*),leftmargin=*]
%%		\item Justifier que $A$, $B$ et $P$ sont alignés.
%%		\item Montrer que $4\vv{PA}+\vv{AB}=\vv{0}$
%%		\item En déduire que $\vv{AP}=\frac{1}{4}\vv{AB}$.
%%	\end{enumerate}
%\end{minipage}
\noindent\parbox{1.0\columnwidth}{
\centering
\begin{tikzpicture}[scale=1.0]
\tkzInit[xmin=-3, xmax=9,ymax=1.5]%\tkzGrid[sub]
\tkzSetUpPoint[size=4, fill=black]
\tkzfctset{tan style/.style={->,>=latex, color=red, line width=1.2pt}}  
\tkzDrawX[left=5pt, label=, font=\scriptsize, line width=1.2pt ] 
\def \xa{4}; 
\def \ya{0};  
\tkzDefPoint(0,0){A} \tkzDefShiftPoint[A](\xa,\ya){B}  
\tkzDrawSegments[-, >=latex, line width=1.2pt](A,B)


\tkzLabelPoints[above, font=\scriptsize](A,B) 

\foreach \i in {0,...,4}{%
	\tkzMarkSegment[pos=\i/4,mark=|,  size=3pt](A,B) 
} 
\end{tikzpicture} 
}

\vfill 
%\vfill 

\end{example}
\begin{exercise}\hfill\\
Pour chaque cas, exprimer $\vv{AP}$ en fonction de $\vv{AB}$, et placer le point $P$ sur la figure.
\begin{multicols}{2}
\begin{enumerate}[label=\alph*),leftmargin=*]
\item $P$ est tel que $ 2\vv{PA}+\vv{PB}=\vv{0}$.
\item $Q$ est tel que $\vv{QA}+2\vv{QB}=\vv{0}$.
\item $R$ est tel que $\vv{RA}+5\vv{RB}=\vv{0}$.
\item $S$ est tel que $-3\vv{SA}+\vv{SB}=\vv{0}$.
\item $T$ est tel que $ -\vv{TA}+3\vv{TB}=\vv{0}$.
\item $M$ est tel que $ \vv{MA}+\vv{MB}=\vv{0}$.
\end{enumerate} 
\end{multicols}
\begin{minipage}{1\columnwidth}
\centering
\begin{tikzpicture}[scale=0.8]
\tkzInit[xmax=6,ymax=1.5]%\tkzGrid[sub]
\tkzSetUpPoint[size=4, fill=black]
\tkzfctset{tan style/.style={->,>=latex, color=red, line width=1.2pt}}  

\def \xa{6}; 
\def \ya{0};  
\tkzDefPoint(0,0){A} \tkzDefShiftPoint[A](\xa,\ya){B} 
\tkzDefPoint(12,0){C}  
\tkzDefPoint(-6,0){D}  
\tkzDrawSegments[-, >=latex, line width=1.2pt](C,D)


\tkzLabelPoints[above, font=\scriptsize](A,B) 

\foreach \i in {0,...,6}{%
	\tkzMarkSegment[pos=\i/6,mark=|,  size=3pt](A,B)
	\tkzMarkSegment[pos=\i/6,mark=|,  size=3pt](A,D)
	\tkzMarkSegment[pos=\i/6,mark=|,  size=3pt](B,C)
} 
\end{tikzpicture} 
\end{minipage}
\end{exercise}
%\vfill
%%%%%%%%%%%%%%%%%%%%%%%%%%%%%%%%%%%%%%%%%%%%%%%%%%%%%%%%%%%
%
%%%%%%%%%%%%%%%%%%%%%%%%%%%%%%%%%%%%%%%%%%%%%%%%%%%%%%%%%%%
%\newpage
\begin{example}[Équation vectorielle dans le plan]\hfill \\
Soit les points $A(1;2)$ , $B(3;1)$ et $C(4;5)$ dans le repère $(O,\vec{\imath},\vec{\jmath})$.\\
On cherche les coordonnées des points $M(x;y)$ et $N(x;y)$ tel que 
\begin{align*}
3\vv{BM} +2\vv{AM}&=\vv{BC} & \vv{BN}+3\vv{NA}&=\vv{BC} 
% 3(\vv{BO}+\vv{OM}) + 2 (\vv{AO}+\vv{OM})& = \vv{BC}\\
% 3\vv{BO}+3\vv{OM} + 2\vv{AO}+2\vv{OM}& = \vv{BC}\\
% 3\vv{OM}+2\vv{OM} + 3\vv{BO}+ 2\vv{AO}& = \vv{BC}\\
%  5\vv{OM} + 3\vv{BO}+ 2\vv{AO}& = \vv{BC}\\
%    5\vv{OM} & = \vv{BC}- 3\vv{BO}- 2\vv{AO}\\
%    5\vv{OM} & = \vv{BC}+ 3\vv{OB}+ 2\vv{OA}\\
%    \vv{OM} & = \frac{1}{5}\left(\vv{BC}+ 3\vv{OB}+ 2\vv{OA}\right)\\
\end{align*}  
\vfill
\end{example}  
%%%%%%%%%%%%%%%%%%%%%%%%%%%%%%%%%%%%%%%%%%%%%%%%%%%%%%%%%%%
%
%%%%%%%%%%%%%%%%%%%%%%%%%%%%%%%%%%%%%%%%%%%%%%%%%%%%%%%%%%%
%\clearpage 
\begin{exercise} \hfill \\
On donne les points $A(-5;2)$. $B(3;0)$ et $C(-1;4)$. Calculer les coordonnées du point $M(x;y)$ qui vérifie l'équation donné.
\begin{multicols}{3} 
\begin{enumerate}[label=\alph*),leftmargin=*]
\item $\vv{CM}=\vv{AB}-3\vv{AC}$
\item  $\vv{MA}+\vv{MB}=\vv{BC}$ 
\item  $\vv{MA}+\vv{MC}=\vv{CB}$
\item $\vv{MA}=3\vv{BM}+\vv{AC}$
\item $2\vv{AM}+\vv{BM}=2\vv{CM}$ 
\item $\vv{MA}+\vv{MB}+\vv{MC}=\vv{0}$
\end{enumerate}
\end{multicols}
\end{exercise}







%----------------------------------------------------------------------------------------
%	 
%----------------------------------------------------------------------------------------

\newpage 
\loadgeometry{widelayout}  
\pagestyle{scrheadings} % Print headers again  
\KOMAoptions{
headwidth=textwithmarginpar,
headsepline=true,
headsepline= 0.5pt,
footwidth=textwithmarginpar,
} 
\onehalfspacing


\section{Club maths : Cas particuliers du théorème de Menelaus}
\begin{exercise}\hfill \\
$ABC$ est un triangle. $P$ est le milieu de $[AB]$, $Q$ est un point de $(AC)$ et $R$ est sur le segment $[BC]$.\\
Les graduations sur les droites sont régulières.\\
L'objectif de cet exercice est de démontrer que les points $P$, $Q$ et $R$ sont alignés.\\
\begin{minipage}{0.5\columnwidth} 
\begin{enumerate}
\item Donner les valeurs des réels $a$, $b$ et $c$ tel que :
\[
\vv{AP}=a\vv{AB}\qquad \vv{AQ}=b\vv{AC}\qquad \vv{BR}=c\vv{BC}
\]
\item Montrer que $\vv{PQ}=-\tfrac{1}{2}\vv{AB}+\tfrac{5}{4}\vv{AC}$.
\item Exprimer $\vv{PR}$ en fonction de $\vv{AB}$ et $\vv{AC}$.
\item Montrer que $\vv{PQ}$ et $\vv{PR}$ sont multiples l'un de l'autre. Conclure.
\end{enumerate}
\end{minipage}
\begin{minipage}{0.5\columnwidth}\centering
\begin{tikzpicture}
\tkzSetUpPoint[size=3, fill=black] 
\tkzDefPoint(0,0){A} 
\tkzDefPoint(1.35,3.25){C}
\tkzDefPoint(5.59,0.47){B} 
\tkzDefPointsBy[homothety=center A ratio 5/4](C){Q}
\tkzDefPointsBy[homothety=center A ratio +1/2](B){P}
\tkzDefPointsBy[homothety=center B ratio +5/6](C){R} 
\tkzDrawPoints(A,B,C) 
\tkzDrawPoints(P,Q,R)  
\tkzLabelPoint[left](A){\scriptsize$A$}
\tkzLabelPoint[right](B){\scriptsize$B$}
\tkzLabelPoint[above left](C){\scriptsize $C$} 
\tkzLabelPoint[below](P){\scriptsize$P$}
\tkzLabelPoint[above](Q){\scriptsize$Q$}
\tkzLabelPoint[above right](R){\scriptsize$R$}
\foreach \i in {1,...,5}{%
	\tkzMarkSegment[pos=\i/6,mark=*,  size=2pt](B,C)
}
\foreach \i in {1,...,5}{%
	\tkzMarkSegment[pos=\i/4,mark=*,  size=2pt](A,C)
}
\tkzMarkSegments[mark=s||,  size=3pt](A,P P,B) 

\tkzDrawSegments[line width=1.2pt](A,Q A,B B,C)
\tkzDrawSegments[line width=0.6pt, dashed](P,Q)
\end{tikzpicture}
\end{minipage}

\end{exercise}  
%%%%%%%%%%%%%%%%%%%%%%%%%%%%%%%%%%%%%%%%%%%%%%%%%%%%%%%%%%%
%
%%%%%%%%%%%%%%%%%%%%%%%%%%%%%%%%%%%%%%%%%%%%%%%%%%%%%%%%%%%
%%%%%%%%%%%%%%%%%%%%%%%%%%%%%%%%%%%%%%%%%%%%%%%%%%%%%%%%%%%
%
%%%%%%%%%%%%%%%%%%%%%%%%%%%%%%%%%%%%%%%%%%%%%%%%%%%%%%%%%%%
%\clearpage 

\begin{exercise}\hfill \\
$ABC$ est un triangle. $P$ est le milieu de $[AB]$, $Q$ est un point de $(AC)$. 
Les graduations sur les droites sont régulières.  
$R$ est le point d'intersection des droites $(PQ)$ et $(BC)$. \\
On sait que pour un certain nombre $k$ : $\vv{BR}=k\vv{BC}$.
L'objectif de cet exercice est de trouver $k$.\\
\begin{minipage}{0.5\columnwidth} 
\begin{enumerate}
\item Donner les valeurs des réels $a$ et $b$ tel que :
\[
\vv{AP}=a\vv{AB}\qquad \vv{AQ}=b\vv{AC} 
\] 
\item Montrer que $\vv{PQ}=-\tfrac{2}{3}\vv{AB}+\tfrac{5}{4}\vv{AC}$.
\item Montrer que $\vv{PR}=\left(\tfrac{1}{3}-k\right)\vv{AB}+k\vv{AC}$.
\item Trouver $k$ pour que les vecteurs $\vv{PQ}$ et $\vv{PR}$ soient multiples l'un de l'autre. 
\end{enumerate}
\end{minipage}
\begin{minipage}{0.5\columnwidth}\centering
\begin{tikzpicture}
\tkzSetUpPoint[size=3, fill=black] 
\tkzDefPoint(0,0){A} 
\tkzDefPoint(1.35,3.25){C}
\tkzDefPoint(5.59,0.47){B} 
\tkzDefPointsBy[homothety=center A ratio 5/4](C){Q}
\tkzDefPointsBy[homothety=center A ratio +2/3](B){P}
\tkzDefPointsBy[homothety=center B ratio +5/7](C){R} 
\tkzDrawPoints(A,B,C) 
\tkzDrawPoints(P,Q,R)  
\tkzLabelPoint[left](A){\scriptsize$A$}
\tkzLabelPoint[right](B){\scriptsize$B$}
\tkzLabelPoint[above left](C){\scriptsize $C$} 
\tkzLabelPoint[below](P){\scriptsize$P$}
\tkzLabelPoint[above](Q){\scriptsize$Q$}
%\tkzLabelPoint[above right](R){\scriptsize$R$}
%\foreach \i in {1,...,6}{%
	%\tkzMarkSegment[pos=\i/7,mark=*,  size=2pt](B,C)
	%}
\foreach \i in {1,...,5}{%
	\tkzMarkSegment[pos=\i/4,mark=*,  size=2pt](A,C)
}
\foreach \i in {1,...,2}{%
	\tkzMarkSegment[pos=\i/3,mark=*,  size=2pt](A,B)
}
%\tkzMarkSegments[mark=s||,  size=3pt](A,P P,B) 

\tkzDrawSegments[line width=1.2pt](A,Q A,B B,C)
\tkzDrawSegments[line width=0.6pt, dashed](P,Q)
\end{tikzpicture}
\end{minipage} 
\end{exercise} 
%%%%%%%%%%%%%%%%%%%%%%%%%%%%%%%%%%%%%%%%%%%%%%%%%%%%%%%%%%%
%
%%%%%%%%%%%%%%%%%%%%%%%%%%%%%%%%%%%%%%%%%%%%%%%%%%%%%%%%%%%
%\clearpage 
\begin{exercise} \hfill\\	
$ABC$ est un triangle ; $P$ est un point de $(AB)$, $Q$ un point de $(BC)$ et $R$ un point de $(AC)$ disposés comme sur le dessin (les graduations sur les droites sont régulières). \\
\begin{minipage}{0.5\columnwidth}
\begin{enumerate}
\item Donner les valeurs des réels $a$, $b$, et $c$ tels que :
\[
\vv{AP}=a\vv{AB}\qquad \vv{AR}=b\vv{AC}\qquad\vv{BQ}=c\vv{BC} 
\]  
\item Exprimer $\vv{PR}$ en fonction de $\vv{AB}$ et $\vv{AC}$.
\item Démontrer que $\vv{PQ}=\frac{9}{28}\vv{AB}+\frac{3}{7}\vv{AC}$
\item Justifier que $\vv{PQ}=-\frac{9}{7}\vv{PR}$. Que conclure ?
\end{enumerate} 
\end{minipage}
\begin{minipage}{0.5\columnwidth}\centering
\begin{tikzpicture}
\tkzSetUpPoint[size=3, fill=black] 
\tkzDefPoint(0,0){A}
%	\tkzDefPoint(1.37,1.33){T}
%	\tkzDefPoint(2.11,0.21){E}
\tkzDefPoint(1.35,3.25){C}
\tkzDefPoint(5.59,0.47){B} 
\tkzDrawPoints(A,B,C)
\tkzLabelPoint[above](C){$C$}
\tkzLabelPoint[left](A){$A$}
\tkzLabelPoint[right](B){$B$}
\tkzDefPointsBy[homothety=center A ratio -1/3](C){R1}
\tkzDefPointsBy[homothety=center A ratio +1/3](C){R2}
\tkzDefPointsBy[homothety=center A ratio +2/3](C){R3}
\tkzDrawPoints(R1,R2,R3)
\tkzLabelPoint[below](R1){$R$}

\foreach \i in {1,2,3,4,5,6}{
	\tkzDefPointsBy[homothety=center B ratio \i/7](C){Q\i}
	\tkzDrawPoint(Q\i)
}
\foreach \i in {1,2,3}{
	\tkzDefPointsBy[homothety=center B ratio \i/4](A){P\i}
	\tkzDrawPoint(P\i)
}
\tkzLabelPoint[below right](P3){$P$}
\tkzLabelPoint[above right](Q3){$Q$}
\tkzDrawSegments[line width=1.2pt](A,B B,C C,R1 P3,Q3 P3,R1)
\end{tikzpicture}
\end{minipage}

\end{exercise}



%----------------------------------------------------------------------------------------
%	 
%----------------------------------------------------------------------------------------
\cleardoublepage 
%----------------------------------------------------------------------------------------
%	
%----------------------------------------------------------------------------------------
\end{document}
